\documentclass[11pt]{amsart}

\usepackage[T1]{fontenc}
\usepackage[utf8]{inputenc}
\usepackage{microtype}
\usepackage{mathtools}
\usepackage{amssymb}
\usepackage{comment}
\usepackage{enumitem}
% Shared review markup and formal links; the human draft hides declaration names.
% Review annotations, following the red-add/blue-remove convention used in
% the Varenna 2016 papers. These commands keep the original wording visible.
% Accept an insertion by removing its AIaddition environment or AIadd wrapper.
% Accept a replacement by deleting AIremove{old} and unwrapping AIadd{new}.
% Resolve a comment by deleting its complete AIcomment{ID}{text} command.
% There is deliberately no global switch that hides unresolved comments.
\usepackage{xcolor}
\usepackage{xurl}
\usepackage[normalem]{ulem}
% Give long inline review markup a little extra line-breaking room without
% changing the author text or the mathematical content.
\setlength{\emergencystretch}{.75em}
% AI compile repair: the retained draft uses the undefined typo \cots.
% EDIT-02 beside that occurrence makes this repair visible in the PDF.
\providecommand{\cots}{\cdots}
% Permit a line break after commas in long inline coordinate lists.
\begingroup\lccode`~=`,
\lowercase{\endgroup\def~}{\mathchar"613B\allowbreak}
\mathcode`,="8000
\definecolor{AIadded}{RGB}{160,25,25}
\definecolor{AIreview}{RGB}{25,70,145}
\newcommand{\AIadd}[1]{{\color{AIadded}#1}}
\newcommand{\AIremove}[1]{{\color{AIreview}%
  \ifmmode\text{\sout{\ensuremath{#1}}}\else\sout{#1}\fi}}
\newcommand{\AIreplace}[2]{\AIremove{#1}\,\AIadd{#2}}
\newenvironment{AIaddition}[1]
  {\par\begingroup\color{AIadded}\noindent
   \textbf{AI addition: #1.}\par\smallskip\ignorespaces}
  {\par\endgroup}
\newcommand{\AIcomment}[2]{%
  \par\smallskip\noindent
  \fcolorbox{AIreview}{AIreview!4}{%
    \parbox{\dimexpr\linewidth-2\fboxsep-2\fboxrule\relax}{%
      \normalfont\small\color{AIreview}\textbf{AI comment #1.} #2}}%
  \par\smallskip}
\newcommand{\AIreviewlegend}{%
  \AIcomment{Review guide}{Red text proposes additions or corrections;
  blue struck text preserves the original wording.
  \textbf{Johanna:} review the proposals as mathematics and prose, and resolve
  their boxes after accepting or revising them. No Lean installation, code
  editing, or proof replay is required for her review.
  \textbf{Max:} review the whole paper--Lean correspondence, including
  every statement, its hypotheses, and the definitions entering the challenge. Comparison
  notes explain proof differences beside the relevant passages; the technical
  supplement gives mathematical accounts and optional declaration links. The
  side-by-side review interface links statements to their checked sources.}}

% Lean cross-references for the paper--Lean audit.
% Each formalized result is followed by \formalref{...}; every \leandecl or
% \leanlib names one declaration by file, line and fully qualified name at
% the pinned revisions below. The review generator in itpplasma/stafford38-formal
% (tools/paper_lean_audit) checks every such reference against the Lean
% source, so a stale file, line or name fails its build.
% Accepting: keep this file and the \formalref lines, or delete both.
\newcommand{\leanrepo}{https://github.com/itpplasma/stafford38-formal/blob/591d4e11e3971577be169e428ac73c44a4ac8a32}
\newcommand{\leanlibrepo}{https://github.com/itpplasma/algebraic-analysis/blob/4aae47967f6ba02ffe2f639ab06564c9a9d1ecc8}
% \leandecl{file}{line}{Fully.Qualified.Name}: declaration in stafford38-formal.
\newcommand{\leandecl}[3]{\href{\leanrepo/#1\#L#2}{\nolinkurl{#3}}}
% \leanlib{file}{line}{Fully.Qualified.Name}: declaration in AlgebraicAnalysis.
\newcommand{\leanlib}[3]{\href{\leanlibrepo/#1\#L#2}{\nolinkurl{#3}}}
% The human-readable paper displays only substantive mathematical qualifications.
% Full file paths and declaration names remain in source and in the supplement.
\newif\ifshowleandetails
\showleandetailsfalse
% \formalref[relation to the printed statement]{declarations}
\newcommand{\formalref}[2][]{%
  \ifshowleandetails
    \par\smallskip\noindent
    {\small\textsf{Lean}\if\relax\detokenize{#1}\relax\else\ (#1)\fi: #2}%
    \par\smallskip
  \else
    \if\relax\detokenize{#1}\relax\else
      \par\smallskip\noindent{\small\textsf{Formal comparison:} #1}%
      \par\smallskip
    \fi
  \fi}

\usepackage[colorlinks=true,citecolor=blue,linkcolor=blue,urlcolor=blue]{hyperref}
\numberwithin{equation}{section}

\newtheorem{theorem}{Theorem}[section]
\newtheorem{proposition}[theorem]{Proposition}
\newtheorem{lemma}[theorem]{Lemma}
\newtheorem{corollary}[theorem]{Corollary}
\theoremstyle{plain}
\newtheorem{definition}[theorem]{Definition}
\newtheorem{convention}[theorem]{Convention}
\theoremstyle{remark}
\newtheorem{remark}[theorem]{Remark}



\title{Stafford's conjecture on cyclicity of torsion modules}
\author{Christopher Albert et al.}
\address{Graz University of Technology}
\date{}
\newcommand{\Char}{\operatorname{Char}}
\newcommand{\Weyl}{A_n}
\newcommand{\ord}{\operatorname{ord}}
\newcommand{\gr}{\operatorname{gr}}
\newcommand{\Ann}{\operatorname{Ann}}
\newcommand{\Dir}{\operatorname{Dir}}
\newcommand{\Sp}{\operatorname{Sp}}
\newcommand{\Spec}{\operatorname{Spec}}
\newcommand{\Supp}{\operatorname{Supp}}
\newcommand{\Sing}{\operatorname{Sing}}
\newcommand{\sm}{\mathrm{sm}}
\newcommand{\bbA}{\mathbb{A}}
\newcommand{\bbP}{\mathbb{P}}
\newcommand{\bbN}{\mathbb{N}}
\newcommand{\bbZ}{\mathbb{Z}}
\newcommand{\bbQ}{\mathbb{Q}}
\newcommand{\Kbar}{\overline{k}}
\newcommand{\dd}{\mathrm{d}}
\DeclareMathOperator{\rad}{rad}


\begin{document}

\begin{abstract}
Let $A_n(k)$ be the Weyl algebra over a field of characteristic zero.  We prove
that every nonzero $d\in A_n(k)$ admits $F,R,S\in A_n(k)$ such that
$1=dR+FdS$.  Thus Stafford's Conjecture 3.8 holds.  \AIadd{For $n\ge1$,} \AIreplace{The}{the} multiplier $F$ may be
chosen as $x^N$, after a linear symplectic change of variables, where $N$ is
the Bernstein degree of $d$.  The proof studies a cyclic quotient through its
characteristic variety and excludes nonzero quotients by conormal geometry at
infinity.
\end{abstract}

\maketitle
\AIreviewlegend
\AIcomment{AUDIT-01}{\textbf{Information --- no edit requested.} Corrections, clarifications and mathematical differences
from the formal proof remain marked as AI additions. Long formal proof
comparisons and declaration links are collected in the separate Overleaf
supplement \texttt{lean\_proof\_details.tex}. The human-readable draft
shows the mathematical qualifications beside the relevant passages; the
side-by-side audit retains the full technical correspondence.}

\section{Introduction}

Let $k$ be a field of characteristic zero, and let
\[ A_n(k)=k\langle x_1,\ldots,x_n,\partial_1,\ldots,\partial_n\rangle\]
be the $n$th Weyl algebra, with $[\partial_i,x_j]=\delta_{ij}$ and
$[\partial_i,\partial_j]=[x_i,x_j]=0$.

\AIremove{Stafford conjectured in 1978 (Conjecture~3.8 of
\cite[p.~438]{Stafford1978}) that for every nonzero $d\in A_n(k)$, one can
find $F\in A_n(k)$ with $A_n(k)=dA_n(k)+FdA_n(k)$. Stafford notes that
this immediately implies $A_n(k)/dA_n(k)$ is cyclic, i.e.\ generated by a
single element. \cite{Bellamy2026} reported the resulting question --
whether every right torsion module over $A_n(k)$ is cyclic -- as open
[Section~3, Conjecture~3.4]. We prove Stafford's conjecture:}\par
\AIadd{Stafford's Conjecture~3.8 asks whether every nonzero $d\in A_n(k)$
admits $F\in A_n(k)$ with $A_n(k)=dA_n(k)+FdA_n(k)$
\cite[p.~438]{Stafford1978}. Stafford identifies this with stable cyclicity
for finitely generated torsion right Weyl-algebra modules
\cite[pp.~437--438]{Stafford1978}. Bellamy's survey
states the weaker ordinary-cyclicity question for finitely generated torsion
$A_n(\mathbb C)$-modules as open~\cite[Section~3, Conjecture~3.4]{Bellamy2026}.
We prove Stafford's conjecture:}
\AIcomment{INTRO-01}{\textbf{Proposed change --- author acceptance pending.} The quotient $A_n/dA_n$ is always cyclic, generated
by the class of $1$. The nontrivial implication concerns general
\emph{finitely generated} torsion modules. Bellamy states the ordinary
cyclicity question over $\mathbb C$; Stafford's stable-cyclicity formulation
is stronger. The red replacement above restores these distinctions.}

\begin{theorem}\label{thm:main}
For every characteristic-zero field $k$, every $n\geq1$, and every
nonzero $d\in A_n(k)$, there exist $F,R,S\in A_n(k)$ such that
\[ 1=dR+FdS.\]
Moreover, one may choose $F=\ell^{\deg_Bd}$ for a linear Weyl coordinate
$\ell$, obtained by an invertible linear symplectic change of generators.
\AIadd{Explicitly, $\ell=\Phi_g^{-1}(x_1)$ for a linear
$\omega$-preserving automorphism $g$ as in Section~\ref{sec:monic}.}
\end{theorem}
\AIadd{\formalref[{Lean also covers $n=0$, where $A_0(k)=k$}]{%
\leandecl{Stafford38/FoundationClosure.lean}{53}{Stafford38.universalStatement} proves
\leandecl{Stafford38/Statement.lean}{19}{Stafford38.UniversalStatement};
\leandecl{Stafford38/FoundationClosure.lean}{57}{Stafford38.universalFixedSourceStatement} proves
\leandecl{Stafford38/FixedSourceStatement.lean}{36}{Stafford38.FixedSource.UniversalFixedSourceStatement}
with \leandecl{Stafford38/FixedSourceStatement.lean}{23}{Stafford38.FixedSource.IsLinearWeylCoordinate}
and \leandecl{Stafford38/FixedSourceStatement.lean}{18}{Stafford38.FixedSource.bernsteinDegree};
Mathlib-only statements compared by Comparator:
\leandecl{Challenge.lean}{63}{Stafford38Challenge.UniversalStatement},
\leandecl{FixedSourceChallenge.lean}{116}{Stafford38FixedSourceChallenge.UniversalFixedSourceStatement}.}}

\AIadd{The machine-checked proof of this result is registered in
Palomar~\cite{Stafford38Palomar2026,Stafford38Zenodo120}. Its Global Stafford extension
\cite{GlobalStaffordPalomar2026,GlobalStaffordZenodo106} proves the same-divisor identity
$1=dR+FdS$ in $D_k(A)$ for every smooth integral finitely generated
characteristic-zero $k$-algebra $A$ and every nonzero finite-order
$k$-linear differential operator $d$. Polynomial $A$ recovers the
Weyl-algebra identity; the global theorem asserts existence of $F$,
without the fixed linear-coordinate refinement in Theorem~\ref{thm:main}.}

\textit{Proof sketch:}
Give each standard generator Bernstein degree one and
let $N=\deg_B d$.  A linear symplectic change of variables makes $d$ monic of
degree $N$ in a momentum $p$, whose conjugate coordinate we denote by $x$.  It is
enough to prove that the cyclic right module
\begin{equation} Q=A_n/(dA_n+x^NdA_n).
 \label{eq:Q-intro}\end{equation}
vanishes.  An Euler identity makes $x$ surjective on $Q$, while monicity makes
$H=\{x=0\}$ noncharacteristic and places $\Char(Q)$ in a hypersurface
$V(P)$ with $P(dx)\ne0$.  A filtered-page length argument shows that
$\Char(Q)$ has no point over $H$.  Gabber's theorem and the conormal
argument of \AIreplace{Sections 4 and 5}{Sections~\ref{sec:conormal-review}
and~\ref{sec:involutive-review}} then force $\Char(Q)$ to be empty.  The same
argument gives the more precise identity $A_n=dA_n+x^NdA_n$ in the chosen
symplectic coordinates.

\AIadd{Constructive predecessors concern Stafford's established generator
theorems. Hillebrand--Schmale and Leykin give effective two-generator
methods, and Leykin also constructs cyclic generators of holonomic
modules~\cite{HillebrandSchmale2001,Leykin2004}. Quadrat--Robertz develop
module-decomposition and basis algorithms~\cite{QuadratRobertz2007,
QuadratRobertz2014}; Caro--Levcovitz and Caro-Tuesta--Levcovitz study
differential-operator rings with series coefficients~\cite{CaroLevcovitz2010,
CaroTuestaLevcovitz2020}. The constrained identity for arbitrary nonzero $d$
requires the additional quotient-vanishing argument developed here.}

\AIadd{The surrounding generator theory includes Stafford's stable-range
and noncommutative Forster--Swan results~\cite{StaffordStableRange1981,
StaffordGenerating1981}. The generator article has a corrigendum and
related exposition~\cite{StaffordGeneratingCorrigendum1983,
StaffordGeneratingChapter1982}. Stafford also constructed stably free
Weyl-algebra ideals~\cite{StaffordStablyFree1985}.
OreModules supplies computational methods for modules over Ore algebras
\cite{OreModules2007}. These provide context; the Euler residue,
filtered-length contradiction, and asymptotic conormal construction below
are project arguments.}

\section{\texorpdfstring{\AIremove{AI Statement [Chris]}\AIadd{Formal verification}}{Formal verification}}
\AIreplace{Every statement of this paper, including
Theorems~\ref{thm:asymptotic-conormal} and \ref{thm:coisotropic-exclusion}
at the generality stated, is formalized and verified in Lean~4 against
Mathlib, with no axiom beyond Lean's three standard ones.}{The main identity, its fixed-source strengthening and the documented
cyclicity and noncharacteristic consequences are verified in Lean~4 against
Mathlib using only the three standard axioms. The formal proof also supplies
the characteristic-support and geometric results used by its own argument.
Statement scope and proof routes must be compared separately: machine
verification of a terminal theorem does not certify every intermediate
step of the printed proof. The mathematical differences are noted below;
full formal comparisons and declaration links are in the technical supplement.}
\AIadd{The headline theorem is registered in Palomar~\cite{Stafford38Palomar2026};
the current Lean sources and auxiliary results are archived~\cite{Stafford38Zenodo120}.}
The accompanying proof-map supplement names the Lean declaration for each step.
\AIcomment{WF-01}{\textbf{Proposed change --- author acceptance pending.} Please review each comparison at the stated generality.
The Palomar record checks the registered main theorem; later results have
separate archived sources. The formal geometric route uses Laurent series,
while the printed tangent-limit argument is auxiliary. The short comparison
notes identify these differences without claiming that every numbered
statement or every proof passage has been verified verbatim.}
\AIadd{The development uses Lean~4~\cite{Lean4}, Mathlib~\cite{Mathlib2020}
and AlgebraicAnalysis v0.3.0~\cite{AlgebraicAnalysis030}; tooling, chronology
and statement comparison are described in Appendix~\ref{app:ai-workflow}.}



















\AIadd{\paragraph{Online resources.} Online Resource~1 is the interactive paper--Lean correspondence review, including the exact challenge definitions and archived proof endpoints; Online Resource~2 is the readable mathematical account of the selected formal proofs. Both are available from the \href{https://itpplasma.github.io/stafford38-formal/}{review companion}, with a \href{https://github.com/itpplasma/stafford38-supplementary}{frozen supplementary package} for journal review. The reusable HTML generator is maintained separately under the Apache~2.0 license; its software archive and the manuscript supplement have distinct version records. Proposed corrections and unresolved review actions remain explicitly marked in this draft.}

\section{Notation and definitions}\label{sec:setting}


Throughout, $k$ is a field of characteristic zero and for the $n$-th Weyl algebra over $k$, we write
\[ A_n(k)=k\langle x_1,\dots,x_n,\partial_1,\dots,\partial_n\rangle,
  \quad [\partial_i,x_j]=\delta_{ij},\quad [x_i,x_j]=[\partial_i,\partial_j]=0 .\]
All ideals are \emph{right} ideals and all modules are \emph{right} modules
unless the contrary is said explicitly. 

\begin{convention}\label{conv:pair}
We single out the first Weyl pair $(x_1, \partial_1)$ and write
\[\qquad p:=\partial_1,\qquad
  C:=A_{n-1}(k)[x_1],\qquad A_n(k)=C[p;\partial_1],\]
so that $x_1$ is central in $C$, $[p,x_1] = 1$, and the transverse algebra $A_{n-1}(k)$ commutes with $x_1$ and $p$. 
Every element of $A_n(k)$ can uniquely be written as $\sum_j p^j c_j$ with $c_j\in C$. 
\AIadd{This is the derivation Ore presentation; Ore\cite{Ore1933} supplies its historical algebraic framework.}
\end{convention}
\AIadd{\formalref[{coefficients on the left, $\sum_j c_jp^j$; both normal forms are unique}]{%
\leandecl{Stafford38/Weyl/IteratedEquivalence.lean}{637}{Stafford38.WeylIteratedEquivalence.presentedIteratedEquiv},
\leandecl{Stafford38/Weyl/OuterOreMonic.lean}{44}{Stafford38.WeylOuterOreMonic.presentedOuterPolynomial}}}

\begin{convention}\label{conv:H}
\AIremove{We define $H:=\{x_1=0\}=\{(0,x_2,\dots,x_n)\}\subseteq\bbA^n$ as a
distinguished hyperplane; its conormal covector is $\dd x_1$, i.e.\
$\xi=(1,0,\dots,0)$. 
We write $\pi\colon T^*\bbA^n\to\bbA^n$
for the projection and identify $\xi_i$ with $\dd x_i$, so that a covector is
$\sum_i\xi_i\,\dd x_i$.}\AIadd{The hyperplane $H=\{x_1=0\}$, the projection
$\pi$ and the identification $\xi_i\leftrightarrow\dd x_i$ are fixed in
Section~\ref{sec:cotangent} below.}
\end{convention}
\AIcomment{DUP-01}{\textbf{Proposed change --- author acceptance pending.} Convention~\ref{conv:H} repeated the definitions of $H$, $\pi$ and
$\xi_i$ given again in Section~\ref{sec:cotangent}; the red text keeps one copy.
Alternatively delete the whole convention environment.}
\subsection{Cotangent space and symplectic structure}\label{sec:cotangent}

Let $X=\mathbb A^n_k$ be affine $n$-space over a field $k$ of characteristic zero, with coordinates
$x_1,\ldots,x_n$. Its \emph{cotangent bundle} $T^*X$ is the space of pairs
$(y,\xi)$ consisting of a point $y\in X$ together with a covector $\xi$
at $y$; since $X$ is affine space, this cotangent space is the same
vector space $(k^n)^*$ at every point, and we may identify
\[
 T^*X \;\cong\; X\times(k^n)^*.
\]
We write $\xi_1,\ldots,\xi_n$ for the coordinates on $(k^n)^*$ dual to
$x_1,\ldots,x_n$, so that a covector is $\sum_i\xi_i\,dx_i$, and identify
$\gr^BA_n(k)=k[x,\xi]$ with the coordinate ring of $T^*X$. We write
$\pi\colon T^*X\to X$ for the projection forgetting the covector, and for
$y\in X$ write $T_y^*X=\pi^{-1}(y)\cong(k^n)^*$ for the \emph{fibre} of
$T^*X$ over $y$.
\AIadd{\formalref{\leandecl{Stafford38/Characteristic/Polynomial.lean}{18}{Stafford38.Characteristic.SymbolRing}
(the ring $k[x,\xi]$; generator \texttt{.inl i} is $x_i$, \texttt{.inr i} is $\xi_i$)}}

We single out the first coordinate $x_1$ and define the distinguished
hyperplane
\[
 H:=\{x_1=0\}=\{(0,x_2,\ldots,x_n)\}\subseteq X;
\]
its conormal covector is $dx_1$, i.e.\ $\xi=(1,0,\ldots,0)$.


\begin{definition}[Filtration {[Def.~2.6 and ~2.9 \cite{SchnellNotes}]}]\label{def:filtration}
A filtration $F_\bullet$ = $A_n(k)F_\bullet$ on $A_n(k)$ is a sequence of linear subspaces 
\[\{0\} = F_{-1} \subseteq F_0 \subseteq F_1\subseteq...\subseteq A_n(k),\]
such that $F_i\cdot F_j \subseteq F_{i+j}$ and $A_n(k) = \bigcup F_j.$
The associated graded algebra, also called symbol ring, is defined as 
\[\mathrm{gr}^F A_n(k) = \bigoplus_{j = 0}^\infty F_j/F_{j-1}.\]
Let $a\neq 0 \in A_n(k)$ and let $j$ be the smallest integer, such that $a \in F_j$, called a \emph{filtration degree} $\deg_F a$. \AIreplace{The the}{The} class of $a$ in $F_j/F_{j-1}$ is called its \emph{symbol} $\sigma_j(a)$.
\end{definition}

In the proof, we use two filtrations of $A_n(k)$.
\begin{itemize}[nosep,leftmargin=1.4em]
\item The \emph{Bernstein filtration} $B_\bullet$ assigns degree $1$ to every
  $x_i$ and every $\partial_i$. The Bernstein degree of $d$ is
  written $\deg_B d$.
\item The \emph{order filtration} $F_\bullet$ assigns degree $0$ to every $x_i$
  and degree $1$ to every $\partial_i$; $F_m$ is spanned by the monomials
  $x^\alpha\partial^\beta$ with $|\beta|\le m$.  %We write $\sigma_m(a)$ for the order-$m$ principal symbol.
\end{itemize}
\AIadd{\formalref[{as spans of ordered PBW words of bounded weight}]{%
\leandecl{Stafford38/Weyl/Filtration.lean}{278}{Stafford38.WeylFiltration.bernsteinPiece},
\leandecl{Stafford38/Weyl/Filtration.lean}{282}{Stafford38.WeylFiltration.orderPiece},
\leandecl{Stafford38/FixedSourceStatement.lean}{18}{Stafford38.FixedSource.bernsteinDegree}}}
For both the order filtration and the Bernstein filtration, the associated graded algebra $\mathrm{gr} A_n(k)$ is simply the polynomial ring $k[x,\xi]=k[x_1,\dots,x_n,\xi_1,\dots,\xi_n]$ in $2n$ variables with $\xi_i$ the symbol of $\partial_i$. In particular, $\mathrm{gr} A_n(k)$ is commutative \cite[\AIadd{Prop.~2.10}]{SchnellNotes}.
\AIadd{Good filtrations and characteristic varieties are treated systematically in Bj\"ork and Hotta--Takeuchi--Tanisaki\cite{Bjork,HTT}.}

\begin{definition}\label{def:good-filtration}
A filtration of a right $A_n$-module $M$ compatible with $F_\bullet$
is an exhaustive increasing chain of \AIreplace{$k_0$}{$k$}-subspaces $M_\bullet$ with
$M_iF_j\subseteq M_{i+j}$\AIadd{ and $M_i=0$ for $i<0$}. It is \emph{good} if the associated graded module
$\gr M=\bigoplus_iM_i/M_{i-1}$ is finitely generated over $\gr^FA_n$. 
\end{definition}
As an example, the filtration $M_i:=qF_i$ of a cyclic module over $A_n$ induced by its generator $q$ is good, since $\gr Q$ is generated by the element $\bar{q}$ denoting the class of $q$ in $Q_0/Q_{-1}.$

For a finitely generated right $A_n$-module $M$ with a good order filtration
$M_\bullet$, the graded module $\gr M$ is a finitely generated \AIreplace{$k_0[x,\xi]$}{$k[x,\xi]$}-module. 
\begin{definition}\label{def:char-variety}
    The \emph{characteristic variety} (also called symbol variety) of a finitely generated $A_n$-module $M$ with a good order filtration
$M_\bullet$ is defined by
\[
  \Char(M):=V\bigl(\Ann_{k[x,\xi]}\gr M\bigr)\subseteq
  T^*\bbA^n=\Spec k[x,\xi],
\]
$\Ann_{k[x,\xi]}\gr M = \{f \in k[x,\xi]: f\cdot m = 0\,\forall m \in \gr M\}$ are the \emph{annihilators} of $\gr M$ over $k[x,\xi]$, and $V$ denotes their \emph{vanishing locus}, i.e. the set of points $(x, \xi)$ for which all annihilator polynomials vanish. This set of points lies in the \emph{cotangent bundle} $T^*\bbA^n$ of the affine $n$-space $\bbA^n.$

\AIremove{$\Char(M)$ is independent of the choice of the good filtration.}
\AIadd{We only use $M=A_n/I$ with the filtration $M_i=qF_i$ induced by the
class $q$ of $1$; then $\Ann\gr M$ is the \emph{order initial ideal} of $I$,
spanned by the order symbols of the elements of $I$. Points are prime ideals;
for a non-closed field this is the correct reading of ``the set of points''.}
\end{definition}
\AIadd{\formalref[{for $M=A_n/I$}]{%
\leandecl{Stafford38/Characteristic/InitialIdeal.lean}{185}{Stafford38.CharacteristicInitialIdeal.orderInitialIdeal},
\leandecl{Stafford38/Characteristic/InitialIdeal.lean}{193}{Stafford38.CharacteristicInitialIdeal.orderCharacteristicSupport},
\leandecl{Stafford38/Characteristic/AssociatedGradedModule.lean}{370}{Stafford38.CharacteristicAssociatedGradedModule.annihilator_orderAssociatedGradedModule}}}


\begin{definition}[Symplectic form and Poisson bracket]\label{def:symplectic-poisson}
A bilinear form $\omega$ on a $k$-vector space is \emph{symplectic} if it
is antisymmetric and nondegenerate.

The standard symplectic form on the vector space $k^{2n}$, with basis
$x_1,\ldots,x_n,\partial_1,\ldots,\partial_n$, is determined by the
relations
\[ \omega(x_i,\partial_j)=\delta_{ij},\qquad
 \omega(x_i,x_j)=\omega(\partial_i,\partial_j)=0,\]
analogously to the Weyl relations. Identifying $T^*X\cong X\times(\mathbb
C^n)^*$ via the coordinates $x_1,\ldots,x_n$ and their duals
$\xi_1,\ldots,\xi_n$, on $T^*X$ this $\omega$ takes the familiar form
\[ \omega=\sum_{i=1}^n dx_i\wedge d\xi_i.\]

The associated \emph{Poisson bracket} on $k[x,\xi]$, the coordinate ring
of $T^*X$, is
\[ \{f,g\}=\sum_{i=1}^n\left(
   \frac{\partial f}{\partial \xi_i}\frac{\partial g}{\partial x_i}
   -\frac{\partial f}{\partial x_i}\frac{\partial g}{\partial \xi_i}
 \right).\]
\end{definition}
\AIadd{\formalref[{Lean's bracket is the negative of this one; involutivity is unaffected}]{%
\leandecl{Stafford38/Characteristic/Polynomial.lean}{27}{Stafford38.Characteristic.poissonBracket},
\leandecl{Stafford38/Characteristic/SymplecticCompletion.lean}{25}{Stafford38.CharacteristicSymplecticCompletion.phasePairing}}}

\begin{definition}[involutive]\label{def:involutive}
A reduced closed subset of $T^*\bbA^n$ is \emph{involutive} if its vanishing ideal is closed under the
Poisson bracket, $\{I(W),I(W)\}\subseteq I(W)$. \AIadd{An ideal $J$ is
\emph{involutive} if $\{J,J\}\subseteq J$. This is weaker than being a Poisson
ideal ($\{J,k[x,\xi]\}\subseteq J$): the ideal $(\xi_1,\dots,\xi_n)$ is involutive
but not Poisson.}
\end{definition}
\AIadd{\formalref{\leandecl{Stafford38/Characteristic/PostScalarExtensionPoisson.lean}{30}{Stafford38.Characteristic.PostScalarExtensionPoisson.IsInvolutive};
negative control \leandecl{Stafford38/Geometry/GeneralCoisotropicSetsTest.lean}{160}{Stafford38.Geometry.GeneralCoisotropicSetsTest.zeroSectionIdealOne_not_isPoisson}}}


\AIcomment{ALG-01}{\textbf{Proposed change --- author acceptance pending.} Gabber applies to the radical characteristic ideal:
put $J=\Ann_{k[x,\xi]}\gr M$ and $I=\sqrt J$; then
$\{I,I\}\subseteq I$. The set $\Char(M)=V(J)$ is not an ideal.
Use $I=\sqrt J$ in the componentwise lemma and its proof below.
Minimal primes of $J$ and $\sqrt J$ coincide, so the localization argument
then has the intended hypotheses.}
\subsection{Gabber's theorem}
\begin{theorem}[Gabber {\cite{Gabber1981}}; algebraic form in
 {\AIremove{\texttt{SinghKumar}}\AIadd{\cite{SinghKumar2014}}}; see also {\cite[\AIreplace{\S2.4}{Theorem~2.3.1 and Appendix~D}]{HTT}}\AIremove{, Bj\"ork}]\label{ithm:gabber}
Let $I$ be any right ideal of $A_n(k)$, with $k$ a field of characteristic zero,
and give $M=A_n(k)/I$ its order quotient filtration. Then\AIremove{, the characteristic variety $\Char(M)$ 
is an involutive ideal, i.e.}
\AIremove{$\{\Char(M), \Char(M)\} \subset \Char(M)$.}\AIadd{ the radical
$\sqrt J$ of $J=\Ann_{k[x,\xi]}\gr M$ is involutive,
$\{\sqrt J,\sqrt J\}\subseteq\sqrt J$; equivalently, $\Char(M)=V(J)$ is involutive
in the sense of Definition~\ref{def:involutive}.}
%$ is a coisotropic subvariety of $T^*\bbA^n$ with respect to the standard symplectic structure $\omega$.
%Equivalently, each of the irreducible components of the characteristic variety (i.e. each minimal prime $\mathfrak p$ of $\operatorname{Ann(gr M)}$) is involutive, meaning $\{f,g\}\in p$ for all $f, g \in \mathfrak p$.
%Put
%$J:=\Ann_{K[x,\xi]}\gr M$. Then the radical $\rad(J)$ is involutive:
%\[ \{\rad(J),\rad(J)\}\subseteq\rad(J).\]
\end{theorem}
\AIadd{\formalref[{proved in Lean, not cited; exactly this scope, over every field of characteristic zero}]{%
\leandecl{Stafford38/Characteristic/GabberGlobalAssembly.lean}{221}{Stafford38.Characteristic.GabberGlobalAssembly.weylAssociatedGradedRadicalInvolutivity},
statement \leandecl{Stafford38/Characteristic/CanonicalGabberInvolutivityInterface.lean}{81}{Stafford38.Characteristic.CanonicalGabberInvolutivityInterface.WeylAssociatedGradedRadicalInvolutivity},
minimal primes \leandecl{Stafford38/Characteristic/GabberGlobalAssembly.lean}{141}{Stafford38.Characteristic.GabberGlobalAssembly.minimalPrime_isInvolutive}}}
\AIcomment{LIT-01}{\textbf{Information --- no edit requested.} Statement check: Gabber~\cite{Gabber1981} proves that
$\sqrt{\Ann\gr M}$ is closed under the Poisson bracket for every finitely
generated module $M$ with a good filtration over a filtered ring whose associated
graded ring is commutative, Noetherian and contains $\mathbb Q$; Singh--Kumar
give the algebraic proof. The red statement above is the special case
$M=A_n(k)/I$ with the quotient filtration, which is all the paper uses and
exactly what the formalization proves (following the Singh--Kumar argument:
Rees ring modulo $T^2$, localization at a minimal prime, Artinian truncation,
trace identity).}
\begin{lemma}[Componentwise involutivity]\label{lem:componentwise-involutive}
\AIreplace{Let $M$ be a finitely generated $A_n(k)$-module with a good filtration, and let
$\Char(M)$ be its characteristic variety, so that by Gabber's theorem $\{\Char(M),\Char(M)\}\subseteq\Char(M)$.
If $\mathfrak p$ is a minimal prime of $\Ann_{k[x,\xi]}\gr M$ (equivalently, a minimal prime over $\Char(M)$),}{Let
$J$ be an ideal of $k[x,\xi]$ whose radical is involutive (by
Theorem~\ref{ithm:gabber}, e.g.\ $J=\Ann\gr M$ for $M=A_n(k)/I$). If
$\mathfrak p$ is a minimal prime of $J$,}
then $\mathfrak p$ is itself involutive:
\[\{\mathfrak p,\mathfrak p\}\subseteq\mathfrak p.\]
\end{lemma}

\begin{proof}
Write $R=k[x,\xi]$ and $I=\AIreplace{\Char(M)}{\sqrt J}$. Localize at $\mathfrak p$. Since $\mathfrak p$ is minimal over $I$, the only prime of $R_{\mathfrak p}$ containing $I_{\mathfrak p}$ is $\mathfrak p_{\mathfrak p}$ itself; as $I$ is already radical, so is its localization, giving
\[I_{\mathfrak p} = \sqrt{I_{\mathfrak p}} = \mathfrak p_{\mathfrak p}.\]
The Poisson bracket extends to $R_{\mathfrak p}$: by the Leibniz rule, $\{a/s,\,b/t\}$ is determined by $\{a,b\}$ and the brackets of $a,b$ with $s,t$, giving a well-defined bracket on $R_\mathfrak p$ agreeing with the one on $R$ under $f\mapsto f/1$; in particular $\{f/1,g/1\}=\{f,g\}/1$ for $f,g\in R$, and the involutivity of $I$ localizes to $\{I_\mathfrak p,I_\mathfrak p\}\subseteq I_\mathfrak p$.

Now take $f,g\in\mathfrak p$. Then $f/1,g/1\in\mathfrak p_\mathfrak p=I_\mathfrak p$, so
\[ \{f,g\}/1 = \{f/1,g/1\} \in I_\mathfrak p = \mathfrak p_\mathfrak p.\]
Thus there is $t\notin\mathfrak p$ with $t\{f,g\}\in\mathfrak p$. Since $\mathfrak p$ is prime and $t\notin\mathfrak p$, this forces $\{f,g\}\in\mathfrak p$.
\end{proof}
\AIadd{\formalref[{any ideal with involutive radical; proof without localization}]{%
\leandecl{Stafford38/Characteristic/RadicalMinimalPrimeInvolutivity.lean}{108}{Stafford38.Characteristic.RadicalMinimalPrimeInvolutivity.minimalPrimes_isInvolutive_of_radical_isInvolutive}}}
\AIcomment{COMP-01}{\textbf{Information --- proof comparison; no edit requested.} The formal argument starts with an involutive radical and proves that each minimal prime is involutive directly, without localization. Thus the hypothesis is radical involutivity, not that the whole ideal is a Poisson ideal. \href{https://itpplasma.github.io/stafford38-formal/lean_proof_details.pdf\#nameddest=comparison-COMP-01}{Mathematical proof (COMP-01, p.~1).}}

\section{Mapping to a monic normal form}\label{sec:monic}
Given a fixed nonzero $d \in A_n(k)$ of Bernstein degree $N$, our first goal is to find a map to replace $d$ by a normal form $\tilde{d}$ that is monic in $p$, i.e. of the form $\tilde d=p^N+\sum_{j<N}p^ja_j$ with coefficients $a_j$ in $C$. \\

If $N=0$, then $d \in B_0 = k$ is a nonzero scalar, hence a unit, and $1 = dR+FdS$ holds trivially with \AIreplace{$R = d^{-1}, F = 0$}{$F=1=\ell^0$, $R=d^{-1}$, $S=0$}. We therefore assume $N\geq1$. 
\AIcomment{ALG-07}{\textbf{Proposed change --- author acceptance pending.} In the scalar case choose $F=1$, $R=d^{-1}$,
$S=0$ to satisfy the stronger assertion $F=\ell^0=1$ too.}

We want to construct a well-behaved automorphism of $\gr^BA_n(k)=k[x,\xi]$
from simpler linear algebra. We proceed in three stages:
define an automorphism of the vector space $\mathbb V=k^{2n}$, extend it to
an automorphism of $A_n(k)$ itself, and observe the induced action this
extension leaves on the associated graded ring.

We can identify $\mathbb V = k^{2n}$ with the Bernstein-degree-one part $B_1/B_0$ of
$A_n(k)$, i.e.\ linear combinations of $x_1,\ldots,x_n,\partial_1,\ldots,
\partial_n$; the coordinates of a vector are then the coefficients of
these generators. By this identification, $\mathbb V$ carries the standard symplectic form $\omega$ (Def. \ref{def:symplectic-poisson}). Let $g$ be a $k$-linear automorphism of $\mathbb V$ that preserves the symplectic form, that is $\omega(gu,gv)=\omega(u,v)$ for all
$u,v\in\mathbb V$.

Every such $g$ extends uniquely to a $k$-algebra automorphism $\Phi_g$ of
$A_n(k)$: substituting each generator $u\in\mathbb V$ by its image $g(u)$
keeps the Weyl relations \AIreplace{$[u,v]=\omega(u,v)$}{$[v,u]=\omega(u,v)$ (recall $[\partial_i,x_i]=1=\omega(x_i,\partial_i)$)} intact exactly when $g$
preserves $\omega$, and this substitution on generators determines
$\Phi_g$ on all of $A_n(k)$.

Finally, since $\Phi_g$ sends each generator to another degree-one
element, it preserves the Bernstein filtration, and therefore induces a
well-defined automorphism of the associated graded ring
$\gr^BA_n(k)=k[x,\xi]$. \AIreplace{Concretely, this induced map sends a symbol (a
polynomial function on $\mathbb V$) $s$ to $s\circ g^{-1}$: functions
transform by precomposition with the inverse, since $\Phi_g$ moves the
underlying points of $\mathbb V$ forward by $g$.}{Concretely, write the generators
as $z_1,\dots,z_{2n}$ and let $G$ be the matrix of $g$, $g(z_i)=\sum_jG_{ji}z_j$.
A symbol is a polynomial $s(z)$, and $\Phi_g$ substitutes $z_i\mapsto\sum_jG_{ji}z_j$.
Reading $s$ as a function of coefficient vectors $c\in k^{2n}$, the symbol of
$\Phi_g(d)$ is therefore $c\mapsto s(G^{\mathsf T}c)$, i.e.\ $s\circ g^{\mathsf T}$
(not $s\circ g^{-1}$). For homogeneous $s$ of degree $N$, the value at the basis
vector $e_p$ is the coefficient of $z_p^N$.}
\AIcomment{ALG-08}{\textbf{Proposed change --- author acceptance pending.} The printed law $s\mapsto s\circ g^{-1}$ is the point
convention $\Phi(z)=z\circ g^{-1}$, not the substitution $\Phi_g(u)=g(u)$ defined
above. Counterexample to the printed proof of Proposition~\ref{prop:monic}:
$n=1$, $d=x$, $v=x+\partial$ (so $s(v)=1$), $g=\begin{psmallmatrix}1&-1\\0&1\end{psmallmatrix}$,
which preserves $\omega$ and maps $v$ to $e_p=\partial$; yet $\Phi_g(x)=x$ has no
$p$-term. The red replacement and the corrected proof below use
$s\circ g^{\mathsf T}$, as the formalization does.}
\AIadd{\formalref{\leandecl{Stafford38/Weyl/Symplectic.lean}{43}{Stafford38.WeylSymplectic.standardSymplecticAlgHom},
\leandecl{Stafford38/Weyl/Symplectic.lean}{88}{Stafford38.WeylSymplectic.standardSymplecticAlgEquivOfInverse}
(with explicit inverse), \leandecl{Stafford38/Weyl/SymbolCompatibility.lean}{411}{Stafford38.WeylSymbolCompatibility.standardSymplecticAlgHom_principal_compatibility}}}

\begin{lemma}\label{lem:sympcompl}
\AIreplace{Let $(\mathbb V,\omega)$ be a symplectic $k$-space of dimension $2n$ and $0\neq v\in \mathbb V$.}{Let $\mathbb V=k^{2n}$ carry the standard symplectic form fixed above, and let $0\neq v\in\mathbb V$.}
\AIreplace{Then $v$ is the first vector of a symplectic basis of $\mathbb V$; equivalently, there
exists a $k$-linear $g$ preserving $\omega$ with $g(v)=e$ for any prescribed nonzero $e$.}{For any prescribed nonzero $e$, there exists a $k$-linear automorphism $g$ preserving $\omega$ with $g(v)=e$.}
\end{lemma}

\begin{proof}
Since $\omega$ is nondegenerate and $v\neq0$, the functional $\omega(v,-)$
is not identically zero, so there is $w'\in\mathbb V$ with
$\omega(v,w')=c\neq0$; setting $w:=w'/c$ gives $\omega(v,w)=1$.

Let $U=\operatorname{span}(v,w)$. For $0\neq av+bw\in U$: if $a\neq0$,
$\omega(av+bw,w)=a\neq0$; if $b\neq0$, $\omega(av+bw,v)=-b\neq0$. Either
way $av+bw$ pairs nontrivially with something in $U$, so $\omega|_U$ is
nondegenerate.

Since $\omega|_U$ is nondegenerate, $U\cap U^\perp=0$, and as
$\dim U+\dim U^\perp=\dim\mathbb V$ for a nondegenerate form, this gives
$\mathbb V=U\oplus U^\perp$. If $u\in U^\perp$ pairs to zero with all of
$U^\perp$, it also pairs to zero with $U$ (as $u\in U^\perp$), hence with
all of $\mathbb V$, forcing $u=0$; so $\omega|_{U^\perp}$ is nondegenerate
too.

By induction on $n$ (trivial for $n=1$, where $U=\mathbb V$), a
symplectic basis $e_2,\ldots,f_n$ of $U^\perp$ exists; setting $e_1:=v$,
$f_1:=w$ gives a symplectic basis of $\mathbb V=U\oplus U^\perp$, since
cross-pairings between $U$ and $U^\perp$ vanish by definition of
$U^\perp$.

For the second formulation: given nonzero $v,e\in\mathbb V$, complete
each to a symplectic basis as above. The linear map $g$ sending one basis
to the other, vector by vector, preserves all pairwise values of $\omega$
on basis vectors, hence by bilinearity preserves $\omega$ on all of
$\mathbb V$; and $g(v)=e$.
\end{proof}
\AIadd{\formalref[{second formulation only; the formal proof composes at most two symplectic transvections}]{%
\leandecl{Stafford38/Characteristic/SymplecticCompletion.lean}{232}{Stafford38.CharacteristicSymplecticCompletion.exists_symplectic_mulVec_eq}}}
\AIcomment{COMP-02}{\textbf{Proposed change --- author acceptance pending.} The red statement specializes to the standard symplectic space used by the paper and retains the transitivity formulation proved by the formal endpoint. The arbitrary-space formulation would additionally require a checked coordinate-identification bridge. The stronger basis sentence remains in blue as an unselected alternative; restoring it requires its own formal correspondence. The formal proof establishes transitivity on nonzero vectors by at most two symplectic transvections. It does not prove the separate symplectic-basis formulation. Reversing the sign of the alternating form leaves its isometry group unchanged. \href{https://itpplasma.github.io/stafford38-formal/lean_proof_details.pdf\#nameddest=comparison-COMP-02}{Mathematical proof (COMP-02, p.~1).}}



\begin{proposition}\label{prop:monic}
Let $0\neq d\in A_n(k)$ with $N:=\deg_B d\ge1$. There exists a $k$-linear $\omega$-preserving automorphism $g$ and
$c\in k^{\times}$ such that
\begin{equation}
  \tilde d:=c^{-1}\Phi_g(d)=p^{\,N}+\sum_{j<N}p^{\,j}a_j,
  \qquad a_j\in C,\quad \deg_B a_j\le N-j,\label{eq:monic-d}
\end{equation}
in the notation of Convention~\ref{conv:pair}.
\end{proposition}

\begin{proof}
Let $s\in k[x,\xi]$ be the Bernstein principal symbol of $d$, a nonzero
homogeneous polynomial of degree $N$ on $\mathbb V$. A characteristic-zero field is
infinite, so every nonzero polynomial has a nonvanishing value: choose $v\in \mathbb V$
with $s(v)\neq0$. \AIreplace{By Lemma~\ref{lem:sympcompl} we can map $v$ to the basis vector $e_1$ in $p$-direction via an automorphism $g$. The Bernstein symbol of $\Phi_g(d)$ is
$s\circ g^{-1}$, whose value at $e_1$ is $s(v)\neq0$.}{By Lemma~\ref{lem:sympcompl} there is an $\omega$-preserving $h$ with $h(e_1)=v$, where $e_1$ is the basis vector in the $p$-direction. Let $g$ be the map with matrix $H^{\mathsf T}$; the transpose of a symplectic matrix is symplectic, so $g$ preserves $\omega$. The Bernstein symbol of $\Phi_g(d)$ is
$s\circ g^{\mathsf T}=s\circ h$, whose value at $e_1$ is $s(v)\neq0$.} Put $c=s(v)$; since $c$ is
a nonzero scalar, it is a central unit of $A_n(k)$, and $\tilde d=c^{-1}\Phi_g(d)$
has its Bernstein symbol taking the value $1$ at $e_1$. Since $e_1$ has coordinate
$1$ in the $p$-direction and $0$ elsewhere, that value is precisely the
coefficient of the monomial $p^{\,N}$ in $\tilde d$. Collecting the remaining
terms by their $p$-degree and writing the coefficients on the right, using the
unique normal form $A_n(k)=\bigoplus_j p^{\,j}C$, gives
$\tilde d=p^{\,N}+\sum_{j<N}p^{\,j}a_j$ with $a_j\in C$. The Bernstein bound
$\deg_B a_j\le N-j$ holds because $\deg_B\tilde d=N$ and $\Phi_g$ preserves
$\deg_B$\AIadd{; no cancellation between different $p$-degrees occurs, because
the normal form $\bigoplus_jp^jC$ is unique}.
\end{proof}
\AIadd{\formalref[{monicity is stated as: $\tilde d\in B_N$ and the PBW coefficient of $\partial_1^N$ is $1$}]{%
\leandecl{Stafford38/Weyl/MonicNormalization.lean}{150}{Stafford38.WeylMonicNormalization.exists_normalized_symplectic_image},
\leandecl{Stafford38/Weyl/MonicNormalization.lean}{173}{Stafford38.WeylMonicNormalization.scalar_or_normalized_symplectic_image},
\leandecl{Stafford38/Weyl/PBWMonicBridge.lean}{88}{Stafford38.WeylPBWMonicBridge.IsPBWMonicAt},
\leandecl{Stafford38/Characteristic/SymplecticCompletion.lean}{282}{Stafford38.CharacteristicSymplecticCompletion.exists_symplectic_chart_matrices}}}
\AIcomment{COMP-03}{\textbf{Information --- proof comparison; no edit requested.} The formal normalization uses an explicit symplectic matrix and its inverse. Its symbol substitution is $s(MX)$; with the convention used here this is $s\circ g^{\mathsf T}$. Monicity means that the $p^N$ PBW coefficient is $1$ and the Bernstein degree is at most $N$, which gives the coefficient bounds above. \href{https://itpplasma.github.io/stafford38-formal/lean_proof_details.pdf\#nameddest=comparison-COMP-03}{Mathematical proof (COMP-03, p.~2).}}
\AIremove{All coefficients involved belong to a finitely generated subfield $k_0\subset k$, namely the subfield
generated by the finitely many coefficients of $d$ in the new generators. If we find $R, S, F \in A_n(k_0)$ sufficing Theorem \ref{thm:main}, the same choices extend to $k$ as well, so proving the identity over $k_0$ is sufficient. We therefore replace $k$ by $k_0$, write $A_n:=A_n(k_0)$, and fix an embedding
$k_0\hookrightarrow\mathbb C$.}
\AIadd{Sections~\ref{sec:monic}--\ref{sec:char} work over $k$ itself; to keep the
printed formulas we write $k_0:=k$ and $A_n:=A_n(k_0)$ there.
Sections~\ref{sec:conormal-review} and~\ref{sec:involutive-review} work over an
algebraically closed field $K$ of characteristic zero; the reader may take
$K=\mathbb C$, and Section~\ref{sec:proof} uses $K=\overline k$. No embedding of
$k$ into $\mathbb C$ is needed.}
\AIcomment{FIELD-01}{\textbf{Proposed change --- author acceptance pending.} The formalization never reduces to a finitely generated
subfield: it runs Sections~\ref{sec:conormal-review}--\ref{sec:involutive-review}
over $\overline k$ and descends the order initial ideal (Section~\ref{sec:proof}).
The red text adopts this route; the struck text is no longer needed.}

\begin{remark}\label{rem:transport}
$\Phi_g$ is a ring automorphism, so it preserves written order term by term.
Consequently an identity $1=\tilde dR+x_1^N\tilde dS$ transports back to
\[
  1=d\bigl(c^{-1}\Phi_g^{-1}(R)\bigr)+F\,d\bigl(c^{-1}\Phi_g^{-1}(S)\bigr),
  \qquad F=\Phi_g^{-1}(x_1)^N,
\]
\AIadd{using that $c^{-1}$ is a central scalar.}
\end{remark}
\AIadd{\formalref{\leandecl{Stafford38/UniversalAssembly.lean}{66}{Stafford38.UniversalAssembly.universalStatement_of_canonicalSupportVanishing},
\leandecl{Stafford38/FixedSourceAssembly.lean}{21}{Stafford38.FixedSource.universalFixedSourceStatement_of_canonicalSupportVanishing}}}

We now use the order filtration on $A_n(k_0)$, where coordinates have
degree zero and derivatives have degree one, to characterize $P =
\sigma_{\mathrm{ord}}(\tilde d)$, the order principal symbol of $\tilde
d$. Recall that a symbol, being the class of an element in a single
graded piece, is automatically \emph{homogeneous} of the corresponding
degree, i.e.\ contains only terms of that total degree. 
\begin{definition}\label{def:monic-xi1}
We say a polynomial $P=\sum_{i=0}^Nc_i\,\xi_1^i\in
k_0[\xi_1,\ldots,\xi_n]$, written as a polynomial in $\xi_1$ with
coefficients $c_i \in k_0[\xi_2,\ldots,\xi_n]$, is \emph{monic of degree $N$ in
$\xi_1$} if $c_N=1$.
\end{definition}
\AIadd{\formalref{\texttt{Polynomial.Monic} after \leandecl{Stafford38/Characteristic/NormalSymbolPolynomial.lean}{24}{Stafford38.Characteristic.NormalSymbolPolynomial.normalSymbolAlgEquiv}}}

\begin{lemma}\label{lem:order-symbol}
Let $\tilde d=p^N+\sum_{j<N}p^{\,j}a_j$ with $a_j\in C$, $\deg_Ba_j\le N-j$,
as in \eqref{eq:monic-d}, and set $N=\deg_B\tilde d$. Then $\tilde d$ has
order $N$, and its order principal symbol is a polynomial in $\xi$ alone,
\begin{equation} \sigma_{\mathrm{ord}}(\tilde d)=P(\xi),\qquad
 P\in k_0[\xi_1,\ldots,\xi_n],\qquad \deg P=N.
 \label{eq:P}\end{equation}
$P$ is monic of degree $N$ in $\xi_1$ and in particular
\begin{equation} P(1,0,\ldots,0)=1.
 \label{eq:P-e1}\end{equation}
\end{lemma}

\begin{proof}
Since $\deg_B\tilde d=N$, every PBW monomial $x^\alpha\partial^\beta$
occurring in $\tilde d$ satisfies $|\alpha|+|\beta|\le N$. A monomial
contributing to the order-$N$ part has $|\beta|=N$, hence $|\alpha|=0$: it
is a pure $\partial^\beta$ with scalar coefficient. Thus $\tilde d$ has
order $N$, and its order-$N$ part is a $k_0$-linear combination of
$\partial^\beta$ with $|\beta|=N$, giving $\sigma_{\mathrm{ord}}(\tilde
d)=P(\xi)$, homogeneous of degree $N$ with coefficients in $k_0$.

The monomial $p^N=\partial_1^N$ occurs in $\tilde d$ with coefficient $1$;
every other order-$N$ monomial comes from some $p^ja_j$ with $j<N$ and
$a_j\in C$ (hence involving no further factor of $p=\partial_1$), so it
carries a positive power of some transverse $\partial_i$, $i>1$. Such a
monomial contributes only to coefficients $c_i$ with $i<N$, never to
$c_N$; hence $c_N$ receives no contribution besides the coefficient of
$p^N$ itself, giving $c_N=1$ identically, i.e.\ $P$ is monic of degree $N$
in $\xi_1$. Evaluating at $(1,0,\ldots,0)$ then gives $P(1,0,\ldots,0)=1$.
\end{proof}
\AIadd{\formalref{\leandecl{Stafford38/Characteristic/InitialIdeal.lean}{313}{Stafford38.CharacteristicInitialIdeal.canonical_orderPrincipalComponent_isFibreOnly},
\leandecl{Stafford38/Geometry/ConormalAxisContradiction.lean}{57}{Stafford38.Geometry.ConormalAxisContradiction.canonical_orderPrincipalComponent_isHomogeneous},
\leandecl{Stafford38/Characteristic/NormalSymbolPolynomial.lean}{92}{Stafford38.Characteristic.NormalSymbolPolynomial.canonicalNormalPolynomial_monic},
\leandecl{Stafford38/Geometry/ConormalAxisContradiction.lean}{126}{Stafford38.Geometry.ConormalAxisContradiction.canonical_orderPrincipalComponent_eval_pureMomentumAxis}}}


\section{$x_1$ as a surjective map}\label{sec:euler}
We define $E=x_1p$ and
\begin{equation} I=\tilde dA_n+x_1^N\tilde dA_n,\qquad Q=A_n/I.
 \label{eq:I-Q}\end{equation}
The goal of this subsection will be to show that $Q = Qx_1$ (Proposition \ref{prop:x-surjective}).
\AIadd{\formalref{\leandecl{Stafford38/Weyl/EulerResidue.lean}{22}{Stafford38.WeylEulerResidue.canonicalRightIdeal}
$=\operatorname{span}\{d,x^Nd\}$ as a right ideal}}

The adjoint action of $E$ defines a grading of $A_n$: %adjoint action = commutator
$x_1$ has Euler degree $1$, $p$ has Euler degree $-1$, and the transverse Weyl algebra
has Euler degree zero. Let $A_{\geq0}$ denote the nonnegative part of $A_n$ in this grading. In PBW form (with $x_i$ on the left of $\partial_i$), $A_{\geq0}$ is generated by the transverse algebra, $E$, and $x_1$. \AIadd{Indeed $A_n$ has the basis $b\,x_1^ap^c$ with $b$ transverse, of Euler degree $a-c$; if $a\ge c$ then $x_1^ap^c=x_1^{a-c}(x_1^cp^c)$ and $x_1^cp^c$ is a polynomial in $E$ by Lemma~\ref{lem:euler}. In particular every element of Euler degree $\ge1$ lies in $A_{\ge0}x_1$, using $Ex_1=x_1(E+1)$ to move $x_1$ to the right.} 
By the commutation relations of the Weyl algebra, it holds that
\begin{align}
    Ex_1=x_1(E+1),\qquad x_1E=(E-1)x_1,\nonumber \\
    Ep=p(E-1), \qquad pE=(E+1)p
    \label{eq:nonnegative-pbw}
\end{align}
Additionally, $x_1$ commutes with the transverse algebra, so\AIadd{, together with $x_1E=(E-1)x_1$,} $x_1$ is normal in $A^{\geq 0}.$
\AIadd{\formalref[{Lean uses the subring generated by the transverse algebra, $x_1$ and $x_1p$, and one-sided normality $x_1r=r'x_1$, which is all that is used}]{%
\leandecl{Stafford38/Weyl/EulerSubring.lean}{22}{Stafford38.WeylEulerSubring.eulerSubring},
\leandecl{Stafford38/Weyl/EulerSubring.lean}{68}{Stafford38.WeylEulerSubring.coordinate_normal_eulerSubring}}}
\AIcomment{ALG-02}{\textbf{Proposed change --- author acceptance pending.} The displayed products below have reversed shifts.
For $E=x_1p$ and $[p,x_1]=1$, the correct identities are
$x_1^mp^m=\prod_{i=0}^{m-1}(E-i)$ and
$p^mx_1^m=\prod_{i=1}^{m}(E+i)$.
Already $m=1$ gives $x_1p=E$, $px_1=E+1$.
Adjust the induction accordingly; the next residue proof already uses
these correct products.}
\begin{lemma}\label{lem:euler}
For every $m\ge0$,
\begin{equation}\label{eq:eulerprod}
  x_1^{m}p^{\,m}=\AIreplace{\prod_{i=1}^{m}(E-i)}{\prod_{i=0}^{m-1}(E-i)}, \qquad p^{\,m}x_1^{m}=\AIreplace{\prod_{i=0}^{m-1}(E+i)}{\prod_{i=1}^{m}(E+i)},
\end{equation}
the factors in each product commuting with one another.
\end{lemma}

\begin{proof}
Since both products in \eqref{eq:eulerprod} are polynomials in the single element $E$, their factors commute.
We induct on $m$; both identities are trivial for $m=0$. Assuming the
first identity for $m$,
\AIremove{$x_1^{m+1}p^{\,m+1}
 =x_1\bigl(\prod_{i=1}^m(E-i)\bigr)p$}
\AIremove{${}=\bigl(\prod_{i=1}^m(E-1-i)\bigr)x_1p$}
\AIremove{${}=\bigl(\prod_{i=2}^{m+1}(E-i)\bigr)(E-1)$}
\AIremove{${}=\prod_{i=1}^{m+1}(E-i)$,}
\AIadd{\[x_1^{m+1}p^{\,m+1}
 =x_1\Bigl(\prod_{i=0}^{m-1}(E-i)\Bigr)p
 =\Bigl(\prod_{i=0}^{m-1}(E-1-i)\Bigr)x_1p
 =\Bigl(\prod_{i=1}^{m}(E-i)\Bigr)E
 =\prod_{i=0}^{m}(E-i),\]}
using $x_1E=(E-1)x_1$ repeatedly and the definition of $E$. The second identity follows
identically from $Ep=p(E-1)$\AIadd{:
\[p^{m+1}x_1^{m+1}=p\Bigl(\prod_{i=1}^m(E+i)\Bigr)x_1=\Bigl(\prod_{i=1}^m(E+1+i)\Bigr)px_1=\Bigl(\prod_{i=2}^{m+1}(E+i)\Bigr)(E+1).\]}
\end{proof}
\AIadd{\formalref{\leandecl{Stafford38/Weyl/EulerProductIdentities.lean}{308}{Stafford38.Weyl.EulerProductIdentities.eval_fallingEulerProduct} and
\leandecl{Stafford38/Weyl/EulerProductIdentities.lean}{317}{Stafford38.Weyl.EulerProductIdentities.eval_risingEulerProduct} (the corrected products),
\leandecl{proofs/weyl_pure_power.lean}{331}{Stafford.exists_eulerPolynomial_x_pow_mul_d_pow}}}

\begin{lemma}\label{lem:euler-residue}
There exists an element $U\in A_{\geq0}$ such that
\begin{equation} 1+Ux_1\in I.
 \label{eq:euler-residue}\end{equation}
\end{lemma}

\begin{proof}
From Lemma \ref{lem:euler} we define
\[p^Nx_1^N=\prod_{r=1}^N(E+r)=:R_+(E),\qquad
 x_1^Np^N=\prod_{r=0}^{N-1}(E-r)=:R_-(E).\]
The roots of $R_+$ and $R_-$ are disjoint (the sets \AIreplace{$\{0,-1,\ldots,1-N\}$
and $\{1,\ldots,N\}$}{$\{-1,\ldots,-N\}$ and $\{0,\ldots,N-1\}$} of integers, embedded in the prime field
$\mathbb Q\subset k_0$), so $\gcd(R_+,R_-)=1$ in $\mathbb Q[E]$. By the
extended Euclidean algorithm, there exist $u,v\in\mathbb Q[E]\subset
k_0[E]$ with
\[ R_+(E)u(E)+R_-(E)v(E)=p^Nx_1^Nu(E)+x_1^Np^Nv(E)=1.\]
We write $\tilde d=p^N+\Sigma$ with $\Sigma = \sum_{j<N}p^{\,j}a_j$. Then
\[(\tilde d\,x_1^N)u(E)+(x_1^N\tilde d)v(E)=1+w\in I,\]
where
\[w=\Sigma x_1^Nu(E)+x_1^N\Sigma v(E).\]
We can write each $a_j$ in $\Sigma$ in powers of $x_1$, with coefficients in the transverse Weyl algebra. A term $bx_1^rp^j$ coming from \AIreplace{$a_jp^j$}{$p^ja_j$ (Euler degree does not depend on the order of the factors)} has $r\geq0$ and $j<N$, \AIreplace{so and}{and} multiplication by $x_1^N$ gives that term Euler degree at least $N-j>0$. $u$ and $v$ only depend on $E$, so they have Euler degree 0. Therefore every term of $w$ has positive Euler degree, and thus $w\in A_{\geq0}x_1$ \AIadd{(by the remark after~\eqref{eq:nonnegative-pbw})} - or equivalently, there exists a $U\in A^{\geq0}$, such that $w=Ux_1$.
\end{proof}
\AIadd{\formalref[{the ring identity is Lean's; the factorization uses left coefficients $c\,x^ap^j$}]{%
\leandecl{Stafford38/Weyl/EulerResidue.lean}{37}{Stafford38.WeylEulerResidue.positiveEulerResidue_identity},
\leandecl{Stafford38/Weyl/EulerRemainder.lean}{254}{Stafford38.WeylEulerRemainder.positiveEulerResidue_eq_eulerSubring_mul_coordinate},
\leandecl{Stafford38/Weyl/EulerRemainder.lean}{402}{Stafford38.WeylEulerRemainder.presentedPositiveEulerResidue}}}
\AIcomment{COMP-04}{\textbf{Information --- proof comparison; no edit requested.} The formal argument replaces the Euler grading by the subring generated by the transverse algebra, $x_1$ and $E=x_1p$. It needs only one-sided normality $x_1r=r\prime x_1$ and left-coefficient normal forms. The corrected Euler product identities give a polynomial annihilator with nonzero constant term; no two-sided normality is asserted. \href{https://itpplasma.github.io/stafford38-formal/lean_proof_details.pdf\#nameddest=comparison-COMP-04}{Mathematical proof (COMP-04, p.~2).}}

\begin{proposition}\label{prop:x-surjective}
Right multiplication by $x_1$ on $Q$ is surjective.
\end{proposition}

\begin{proof}
Let $q$ be the class of $1$ in $Q$. By Equation \eqref{eq:euler-residue}, we have
$q=-qUx_1$, so $q\in(qA_{\geq0})x_1$. We use the normality of $x_1$ to show
\begin{equation} qA_{\geq0}=(qA_{\geq0})x_1.
 \label{eq:normal-surj}\end{equation}
The inclusion $\supseteq$ holds because $x_1\in A^{\geq0}$ and
$A^{\geq0}$ is a ring. For $\subseteq$, let $r\in A^{\geq0}$ and choose
$r'$ with $x_1r=r'x_1$ (possible since $x_1$ is normal) then $qr=-(qU)x_1r=-(qUr')x_1\in(qA^{\geq0})x_1$ since $Ur'\in A^{\geq0}$.

$qA_{\geq0}$ is stable under right multiplication by $p$: If $m=m'x_1$, then
\[ mp=m'(x_1p)=m'E\in qA_{\geq0}.\]
It is of course also stable under $C$, since \AIreplace{elements of $C$ have Euler degree $0$}{$C=A_{n-1}(k)[x_1]\subseteq A_{\ge0}$}. Since $C,p$ generate $A_n$, we have $qA_{\geq0}=qA_n=Q$. Equation \eqref{eq:normal-surj} is therefore $Q=Qx_1$.
\AIcomment{ALG-03}{\textbf{Proposed change --- author acceptance pending.} $C=A_{n-1}(k)[x_1]$ is not entirely of Euler degree
zero: $x_1$ has degree $1$. The needed stability still holds because
both the transverse algebra and $x_1$ belong to $A_{\geq0}$,
so $C\subseteq A_{\geq0}$. Replace the degree-zero explanation above.}
\end{proof}
\AIadd{\formalref{\leandecl{Stafford38/Weyl/QuotientTransport.lean}{236}{Stafford38.WeylQuotientTransport.presentedCanonicalRightQuotient_rightMul_coordinate_surjective},
\leandecl{Stafford38/Quotient/EulerSurjectivity.lean}{80}{Stafford38.EulerSurjectivity.rightMul_surjective_of_euler_normal};
generation by $C$ and $p$: \leandecl{Stafford38/Weyl/EulerRemainder.lean}{307}{Stafford38.WeylEulerRemainder.closure_pairOldSubring_union_coordinate_momentum_eq_top}}}
\AIcomment{COMP-05}{\textbf{Information --- proof comparison; no edit requested.} The formal proof works element by element with the transverse--Euler subring. The residue identity gives $1+Ux_1\in I$ for a suitable $U$ in that subring; one-sided normality and the monic relation then give surjectivity of right multiplication by $x_1$. This avoids treating a quotient by a vector subspace as an algebra quotient. \href{https://itpplasma.github.io/stafford38-formal/lean_proof_details.pdf\#nameddest=comparison-COMP-05}{Mathematical proof (COMP-05, p.~3).}}

\section{Vanishing cokernel}\label{sec:cokernel}
\AIadd{In this section and the next, $C$ denotes the commutative symbol ring
$k_0[x,\xi]=T[x_1,\xi_1]$ (not the Ore coefficient algebra of
Convention~\ref{conv:pair}), and Corollary~\ref{cor:charP} of Section~\ref{sec:char},
whose proof uses only Lemma~\ref{lem:order-symbol}, is used freely.}
\AIcomment{NOTATION-01}{\textbf{Proposed change --- author acceptance pending.} $C$ below denotes the commutative symbol ring
$T[x_1,\xi_1]$, whereas Convention~\ref{conv:pair} used $C$ for the
noncommutative coefficient algebra $A_{n-1}(k)[x_1]$. Please rename one
or explicitly announce the change of notation.}
Put $M=\operatorname{gr}Q$ and
\[
 \begin{gathered}
 T=k_0[x_2,\ldots,x_n,\xi_2,\ldots,\xi_n],\qquad C=T[x_1,\xi_1],\\
 V=\ker(R_{x_1}:M\to M),\qquad U=\operatorname{coker}(R_{x_1}:M\to M) =M/x_1M.
 \end{gathered}
\]
With $R_{x_1}$ being the right multiplication \AIreplace{with}{by} ${x_1}$.
If $n=1$, we define $T=k_0$. We want to prove that $U=0$ by comparing
two lengths over a localization of $T$. 
\begin{definition}[\AIreplace{\mbox{\cite{coutinho, cohn}}}{\cite[Tags~02LY, 00J3, 00IV]{Stacks}}]\label{def:length}
    The length $\ell_T$ of a \AIreplace{finite T-module}{$T$-module of finite length} $M$ is defined as the number of steps in a composition series of M such that 
    \[0 = M_0\subsetneq M_1\subsetneq \dots \subsetneq M_{\ell_T} = M\]
    where each successive quotient $M_i/M_{i-1}$ is simple (i.e. has no proper nontrivial submodules). The length is well-defined, as it does not depend on the choice of the composition series. Lengths are additive for short exact sequences.
\end{definition}
\AIadd{\formalref{\texttt{Module.length} (Mathlib); finite length of the localized modules:
\leanlib{AlgebraicAnalysis/Module/MinimalSupportKernelCokernelLengths.lean}{45}{AlgebraicAnalysis.localized_kernel_and_cokernel_isFiniteLength}}}

\begin{lemma}\label{lem:tangential-finiteness}
The modules $U$ and $V$ are finite over $T$, and
$\operatorname{Supp}_T V\subseteq\operatorname{Supp}_T U$.
If $\mathfrak q$ is a minimal prime in $\operatorname{Supp}_T U$, then
$U_{\mathfrak q}$ and $V_{\mathfrak q}$ have finite length over
$T_{\mathfrak q}$.
\end{lemma}
\begin{proof}
The good quotient filtration makes $M = \operatorname{gr} Q$ finitely generated over $C = \operatorname{gr} A_n = k_0[x,\xi] = T[{x_1}, \xi_1]$. 
Since $P$ annihilates $M$ (Corollary \ref{cor:charP}), this action factors
through $C/(P)$, which, because $P$ is monic of degree $N$ in $\xi_1$, is generated over $T[{x_1}]$ by
$1,\xi_1,\ldots,\xi_1^{N-1}$. Thus $M$ is finite over $T[{x_1}]$.
Because $M$ is finite, and $T[{x_1}]$ is a Noetherian ring, the kernel $V$ and cokernel $U$ of multiplication in $M$ by ${x_1}$ are also finite over $T[{x_1}]$. Since they are annihilated by ${x_1}$, they are finite over $T$. %Noetherian theorem 3.5 https://www.cambridge.org/core/services/aop-cambridge-core/content/view/5DD13688B6CC1628C3FD7A318CD479FE/9781139171762c1_p1-19_CBO.pdf/commutative-rings-and-modules.pdf

Suppose ${\mathfrak p}$ is a prime of $T$ where $U_{\mathfrak p} = (M/{x_1}M)_{\mathfrak p} =(M_{\mathfrak p}/{x_1}M_{\mathfrak p}) = 0$. Consequently, $M_{\mathfrak p} = {x_1}M_{\mathfrak p}$, i.e. multiplication by ${x_1}$ is a
surjective endomorphism of the finite module $M_{\mathfrak p}$ over the corresponding
localization of $C$. By \AIreplace{Vasconcelo's theorem \cite{matsumura}[Theorem 2.4]}{Vasconcelos' theorem \cite[Theorem~2.4]{matsumura}}, this endomorphism is also
injective, which means $V_{\mathfrak p} = \ker({x_1}:M_{\mathfrak p}\rightarrow M_{\mathfrak p}) = 0$, i.e. ${\mathfrak p} \notin \operatorname{Supp}_T(V)$, and therefore $\operatorname{Supp}_T V\subseteq\operatorname{Supp}_T U$. 

After localization at a minimal support prime $\mathfrak q$ of $U$, both finite
$T$-modules have support contained in the closed point. They therefore have
finite length.
\end{proof}
\AIadd{\formalref{\leandecl{Stafford38/Characteristic/CanonicalOldTangentialFiniteness.lean}{70}{Stafford38.Characteristic.CanonicalOldTangentialFiniteness.canonical_finite_old_coordinate_kernel_cokernel},
\leandecl{Stafford38/Characteristic/CanonicalTangentialSymbolFiniteness.lean}{113}{Stafford38.Characteristic.CanonicalTangentialSymbolFiniteness.canonical_orderAssociatedGradedModule_finite_tangential_coordinate}}}

\AIcomment{ALG-04}{\textbf{Action required --- unresolved proof or drafting issue.} The printed tangential-finiteness argument needs its finite-$T$-module bridge for $U$ and $V$. State why the monic annihilating symbol makes these modules finite over $T$. The later pages also need their $T$-action and $T$-linear maps; localize those pages, not $Q$. Comparison COMP-06 explains the page action and exhaustion, but does not replace the missing justification in Lemma~\ref{lem:tangential-finiteness}.}
\begin{lemma}\label{lem:page-inequality}
For every minimal prime $\mathfrak q$ in $\operatorname{Supp}_T U$,
\[
 \ell_{T_{\mathfrak q}}(U_{\mathfrak q})
 \le \ell_{T_{\mathfrak q}}(V_{\mathfrak q}).
\]
\end{lemma}

\begin{proof}
To prove the inequality we set up a spectral sequence measuring how $\ker(R_{x_1})$ and
$\operatorname{coker}(R_{x_1})$ compare, order by order, to $V=\ker(x:M\to M)$ and
$U=M/xM$.

Consider the filtered complex $Q\xrightarrow{R_{x_1}}Q$ (right multiplication by ${x_1}$,
between two copies of $Q$), with decreasing filtration $G^p=F_{-p}Q$, where $F$ is
the order filtration. Since ${x_1}$ has order $0$, $R_{x_1}(G^p)\subseteq G^p$ for every
$p$; the induced map on $G^p/G^{p+1}$, summed over $p$, is exactly the action of
${x_1}$ on $M=\operatorname{gr}Q$. This is the crudest approximation to $\ker(R_{x_1})$ and
$\operatorname{coker}(R_{x_1})$: it correctly identifies $V,U$ as approximations, but it
does not yet see elements that vanish to higher order than this naive count
predicts. The spectral sequence refines this approximation in successive stages,
called \emph{pages}: for each $r\ge0$, the source and target pages are
\[
 A_r=\bigoplus_{p\in\mathbb Z}\frac{G^p\cap R_{x_1}^{-1}(G^{p+r})}{G^{p+1}\cap R_{x_1}^{-1}(G^{p+r})},
 \qquad
 B_r=\bigoplus_{p\in\mathbb Z}\frac{G^p}{(G^p\cap R_{x_1}(G^{p-r+1}))+G^{p+1}},
\]
where the degree-$p$ summand of $A_r$ isolates elements of $G^p$ whose image under
$R_{x_1}$ drops $r$ extra filtration levels (candidates for the kernel, accurate to
order $r$), and the degree-$p$ summand of $B_r$ isolates the part of $G^p$ not yet
hit by $R_{x_1}$ from $r-1$ levels shallower.

Applying $R_{x_1}$ to a representative gives a well-defined map $d_r:A_r\to B_r$,
shifting the summand index by $r$; one checks directly from these definitions that
\[
 A_0=B_0=M,\qquad d_0={x_1},\qquad A_{r+1}=\ker d_r,\qquad B_{r+1}=\operatorname{coker}d_r,
\]
so that each page is obtained from the previous one by taking kernel and cokernel
of $d_r$, exactly as in the construction of a spectral sequence from a filtered
complex.


Now we fix a minimal prime $\mathfrak q$ in $\operatorname{Supp}_TU$ and localize at
$T\setminus\mathfrak q$. Since localization is exact, it commutes with kernels,
cokernels, intersections, and preimages, and hence with every operation used to
define $A_r,B_r$ and $d_r$ from $G^p$ and $R_{x_1}$. Consequently the localized objects
$(A_r)_{\mathfrak q},(B_r)_{\mathfrak q}$ are again the pages of the spectral
sequence built the same way from the localized complex
$Q_{\mathfrak q}\xrightarrow{R_{x_1}}Q_{\mathfrak q}$. By
Lemma~\ref{lem:tangential-finiteness}, $A_1=V_{\mathfrak q}$ and
$B_1=U_{\mathfrak q}$ have finite length over $T_{\mathfrak q}$; since each
$A_{r+1}\subseteq A_r$ and $B_{r+1}$ is a quotient of $B_r$, all later pages
have finite length as well.

From here on, all lengths are taken over $T_{\mathfrak q}$, and we write $A_r,B_r,U,V$
for the localized objects $(A_r)_{\mathfrak q},(B_r)_{\mathfrak q},U_{\mathfrak q},V_{\mathfrak q}$.

The maps $B_1\to B_r$ are surjective by construction.
Let $K_r:=\ker(B_1\to B_r)$. Since $B_1\to B_{r+1}$ factors through $B_1\to B_r\to B_{r+1}$, we have
$K_r\subseteq K_{r+1}$, so $(K_r)_{r\ge1}$ is an increasing chain of submodules of $B_1$.
This chain exhausts $B_1$: given $u\in B_1=U$, choose a representative in $Q_{\mathfrak q}$. Since
localization is exact, the surjection $Q=Q{x_1}$ of Proposition~\ref{prop:x-surjective} localizes to
$Q_{\mathfrak q}=Q_{\mathfrak q}{x_1}$, so this representative has an ${x_1}$-preimage in $Q_{\mathfrak q}$; the
same holds for that preimage, and so on. Since the order filtration on $Q_{\mathfrak q}$ (induced from
the order filtration on $Q$) is exhaustive, any fixed element lies in some finite filtration step, and
this iterated-preimage process can be carried out only finitely many times before the resulting class
becomes zero in the corresponding page. Concretely, this shows that $u$ maps to $0$ in $B_r$ for some
$r$ depending on $u$, i.e.\ $u\in K_r$. As $u\in B_1$ was arbitrary, $\bigcup_{r\ge1}K_r=B_1$.
Since $B_1$ has finite length, it is a Noetherian $T_{\mathfrak q}$-module. By the ascending
chain condition, the increasing chain $K_1\subseteq K_2\subseteq\cdots$ stabilizes: there exists a finite $R$ with
$K_R=K_{R+i}$ for all $i\geq 0$. Combined with $\bigcup_rK_r=B_1$, stabilization forces $K_R=B_1$.
Finally, $K_R=B_1$ means $B_1\to B_R$ must be the zero map: by the first isomorphism theorem,
$B_1/K_R\cong\operatorname{im}(B_1\to B_R)$, and $K_R=B_1$ gives $B_1/K_R=0$, so
$\operatorname{im}(B_1\to B_R)=0$. But $B_1\to B_R$ is also surjective, so $B_R=0$.\\

At each stage $r\ge1$ the differential $d_r$ gives an exact sequence
\[ 0\to A_{r+1}\to A_r\xrightarrow{d_r}B_r\to B_{r+1}\to0,\]
which splits at $\operatorname{im}(d_r)$ into two short exact sequences.
Additivity of length gives
$\ell(A_r)=\ell(A_{r+1})+\ell(\operatorname{im}d_r)$ and
$\ell(B_r)=\ell(\operatorname{im}d_r)+\ell(B_{r+1})$; eliminating
$\ell(\operatorname{im}d_r)$ yields, for every $r\ge1$,
\begin{equation}\label{eq:additivity}
 \ell(B_r)+\ell(A_{r+1})=\ell(A_r)+\ell(B_{r+1}).
\end{equation}

Summing this identity over $r=1,\ldots,R-1$ and reindexing the sums $\sum_{r=1}^{R-1}\ell(A_{r+1})=\sum_{r=2}^{R}\ell(A_r)$,
$\sum_{r=1}^{R-1}\ell(B_{r+1})=\sum_{r=2}^{R}\ell(B_r)$, the terms $\ell(A_r),\ell(B_r)$ for $2\le r\le R-1$
cancel on both sides, leaving
\[\ell(B_1)+\ell(A_R)=\ell(A_1)+\ell(B_R).\]
Since $B_R=0$ and $A_1=V$, $B_1=U$, this reads
\[ \ell(U)+\ell(A_R)=\ell(V).\]
As $\ell(A_R)\ge0$, we conclude
\[ \ell_{T_{\mathfrak q}}(U_{\mathfrak q})\le\ell_{T_{\mathfrak q}}(V_{\mathfrak q}).\]
\end{proof}
\AIcomment{ALG-09}{\textbf{Action required --- unresolved proof or drafting issue.} Beyond ALG-04: the exhaustion sentence (``this iterated-preimage
process can be carried out only finitely many times'') is not a proof. One
preimage suffices, taken \emph{before} localizing: if $z\in G^p$ and $z=wx_1$
with $w\in G^{p-r+1}$ (possible for $r$ large, since the filtration is
exhaustive), then the class of $z$ vanishes in $B_r$. See comparison COMP-06. Replace the incomplete exhaustion sentence in the printed proof; this comment alone is not an accepted proof edit.}
\AIadd{\formalref{\leandecl{Stafford38/Characteristic/CanonicalPageEulerInequality.lean}{22}{Stafford38.Characteristic.CanonicalPageEulerInequality.canonicalPage_length_target_le_source};
exhaustion \leanlib{AlgebraicAnalysis/Module/UniformBoundaryVanishing.lean}{103}{AlgebraicAnalysis.exists_uniform_subsingleton_localized}}}
\AIcomment{COMP-06}{\textbf{Action required --- unresolved proof or drafting issue.} Localize only the $T$-modules on the pages, not $Q$. Before localization, choose one preimage $y=zx_1$ and use exhaustiveness to put $z$ in some filtration piece; then $y$ dies on a sufficiently late page. After localization the increasing kernels stabilize because $U_{\mathfrak q}$ has finite length. Telescoping the exact page sequences gives the stated inequality. This supplies the missing exhaustion argument identified in ALG-09. \href{https://itpplasma.github.io/stafford38-formal/lean_proof_details.pdf\#nameddest=comparison-COMP-06}{Mathematical proof (COMP-06, p.~4).}}

On the other hand, we can also establish a different inequality.

\begin{lemma}\label{lem:strict-support-inequality}
For every minimal prime $\mathfrak q$ in $\operatorname{Supp}_T U$,
\[
 \ell_{T_{\mathfrak q}}(U_{\mathfrak q})
 > \ell_{T_{\mathfrak q}}(V_{\mathfrak q}).
\]
\end{lemma}
\begin{proof}
By Gabber involutivity (Theorem \ref{ithm:gabber}) and Lemma \ref{lem:componentwise-involutive}, 
every minimal prime $\mathfrak p$ of $\operatorname{Ann}_C M$ is involutive. It follows, that no $\mathfrak p$ contains ${x_1}$: Assume ${x_1}\in \mathfrak p$. Since $P \in \mathfrak p$,  also the Poisson brackets
\[ \{{x_1}, P\}, \quad \{{x_1}, \{x_1, P\}\}, \dots\]
lie in $\mathfrak p$. $\{{x_1},-\}$ differentiates once with respect to $\xi_1$, and because $P$ is homogeneous of degree $N$ and monic in $\xi_1$, its $N$th derivative is the nonzero constant $N!$ - therefore the $N$th bracket would put \AIreplace{$N!$}{$(-1)^NN!$} in $\mathfrak p$. Since $\mathfrak p$ is a proper prime ideal, no nonzero constant can be in $\mathfrak p$, so we must conclude ${x_1} \notin \mathfrak p$.
This result holds under localization at $T\setminus\mathfrak q$, writing $M'=M_{\mathfrak q}$ for the finite module over $C'=(T\setminus\mathfrak q)^{-1}C$. 

\AIremove{Consider the increasing chain of kernels
$\operatorname{ker}({x_1}:M'\rightarrow M')$ $\subseteq \operatorname{ker}({x_1}^2:M'\rightarrow M')$ $\subseteq \dots$.
Since $M'$ is Noetherian, this chain stabilizes at some point $m$ where $K:= \operatorname{ker}({x_1}^m)$ $= \operatorname{ker}({x_1}^{m+1})=\cots$.}
\AIcomment{EDIT-02}{\textbf{Proposed change --- author acceptance pending.} The undefined command \texttt{\string\cots}
is aliased to \texttt{\string\cdots} in the annotation preamble so this
draft compiles. The kernel-chain argument is repeated immediately below;
keep the fuller second version when resolving this comment.}



Write $M'=M_{\mathfrak q}$, a finite module over $C'=(T\setminus\mathfrak q)^{-1}C$;
this is the same localization used throughout. As shown above, no minimal prime of
$\operatorname{Ann}_C(M)$ contains ${x_1}$, and this persists after localizing at
$T\setminus\mathfrak q$: no minimal prime of $\operatorname{Ann}_{C'}(M')$ contains ${x_1}$
either\AIadd{, since the annihilator of the finite module $M$ commutes with localization}.

The chain $\ker({x_1})\subseteq\ker({x_1}^2)\subseteq\cdots$ of
submodules of $M'$ stabilizes, since $M'$ is Noetherian over $C'$; let
$K:=\ker({x_1}^m:M'\to M')$ at a stage where it stabilizes, so ${x_1}K\subseteq K$ and
$\ker({x_1}^{m+1})=\ker({x_1}^m)=K$. Each quotient $\ker({x_1}^{i+1})/\ker({x_1}^i)$ maps to
$V_{\mathfrak q}=\ker({x_1}:M'\to M')$ injectively via $v\mapsto {x_1}^iv$, so, as $V_{\mathfrak q}$ has
finite length, so does each quotient and hence $K$ itself.

Two elementary facts about $K$:
\begin{itemize}
\item ${x_1}^{-1}K=K$: if $v\in K$ then ${x_1}v\in K$ (as ${x_1}K\subseteq K$); if ${x_1}v\in K=\ker({x_1}^m)$
then ${x_1}^{m+1}v=0$, so $v\in\ker({x_1}^{m+1})=K$.
\item $K\cap {x_1}M'={x_1}K$: "$\supseteq$" is trivial. For "$\subseteq$", if $w={x_1}v\in K$ for
some $v\in M'$, then $v\in {x_1}^{-1}K=K$, so $w\in {x_1}K$.
\end{itemize}

For any endomorphism $T$ of a finite-length module $N$,
additivity along $0\to\ker T\to N\to TN\to0$ and $0\to TN\to N\to N/TN\to0$ gives
$\ell(\ker T)=\ell(N/TN)$. Applying this to $T={x_1}|_K$: since $\ker({x_1})\subseteq K$, we
have $\ker({x_1}|_K)=\ker({x_1})=V_{\mathfrak q}$, so
\[\ell(K/{x_1}K)=\ell(V_{\mathfrak q}).\]

The natural surjection $M'/{x_1}M'\to(M'/K)/{x_1}(M'/K)$ has
kernel $(K+{x_1}M')/{x_1}M'\cong K/(K\cap {x_1}M')=K/{x_1}K$, giving the exact sequence
\[0\longrightarrow K/{x_1}K\longrightarrow U_{\mathfrak q}\longrightarrow
 (M'/K)/{x_1}(M'/K)\longrightarrow0,\]
so by additivity and using $\ell(K/{x_1}K)=\ell(V_{\mathfrak q})$,
\[\ell(U_{\mathfrak q}) = \ell(V_{\mathfrak q}) + \ell\bigl((M'/K)/{x_1}(M'/K)\bigr).\]
It remains to show the last term is positive.

Since $\mathfrak q\in\operatorname{Supp}_TU$,
$U_{\mathfrak q}=M'/{x_1}M'\ne0$; choose $\mathfrak P\in\operatorname{Supp}_{C'}(M'/{x_1}M')$
and a minimal prime $\mathfrak p\subseteq\mathfrak P$ of $M'$. We have shown above that ${x_1}\notin\mathfrak p$, so ${x_1}$ is a unit in the localization at $\mathfrak
p$; as $K$ is annihilated by a power of ${x_1}$, $K_{\mathfrak p}=0$. From
$0\to K_{\mathfrak p}\to M'_{\mathfrak p}\to(M'/K)_{\mathfrak p}\to0$ we get
$(M'/K)_{\mathfrak p}\cong M'_{\mathfrak p}\ne0$, so $\mathfrak p\in
\operatorname{Supp}(M'/K)$; since support is closed under passing to larger primes and
$\mathfrak p\subseteq\mathfrak P$, also $\mathfrak P\in\operatorname{Supp}(M'/K)$, i.e.
$(M'/K)_{\mathfrak P}\ne0$.

Since $\mathfrak P\in\operatorname{Supp}_{C'}(M'/{x_1}M')$ and $M'/{x_1}M'$ is annihilated by
${x_1}$, we have ${x_1}\in\mathfrak P$. If ${x_1}(M'/K)_{\mathfrak P}=(M'/K)_{\mathfrak P}$, then as
${x_1}$ lies in the maximal ideal of the local ring $C'_{\mathfrak P}$, Nakayama's lemma \cite{nakayama}\AIadd{ (see \cite[Tag~00DV]{Stacks})}
would force $(M'/K)_{\mathfrak P}=0$, a contradiction. Hence
$\bigl((M'/K)/{x_1}(M'/K)\bigr)_{\mathfrak P}\ne0$, so $(M'/K)/{x_1}(M'/K)\ne0$.

Therefore $\ell(U_{\mathfrak q})=\ell(V_{\mathfrak q})+\ell\bigl((M'/K)/{x_1}(M'/K)\bigr)
>\ell(V_{\mathfrak q})$.
\end{proof}
\AIadd{\formalref{\leanlib{AlgebraicAnalysis/Module/BaseLocalizedKoszulPositivity.lean}{22}{AlgebraicAnalysis.BaseLocalizedKoszulPositivity.localized_length_cokernel_gt_kernel};
$x_1$ in no minimal prime: \leandecl{Stafford38/Characteristic/NoncharacteristicMinimalPrime.lean}{107}{Stafford38.Characteristic.NoncharacteristicMinimalPrime.canonical_minimalPrime_mem_of_normalCoordinate_false}}}
\AIcomment{COMP-07}{\textbf{Information --- proof comparison; no edit requested.} The formal proof uses Gabber directly for minimal primes. If such a prime contained $x_1$, the $N$-fold bracket of the monic symbol $P$ with $x_1$ would put the unit $\pm N!$ in that prime. The subsequent stable-kernel length computation and Nakayama argument agree with the printed proof. \href{https://itpplasma.github.io/stafford38-formal/lean_proof_details.pdf\#nameddest=comparison-COMP-07}{Mathematical proof (COMP-07, p.~4).}}

\begin{lemma}\label{lem:Uzero}
The cokernel $U=M/x_1M$ of right multiplication by $x_1$ on
$M=\operatorname{gr}Q$ vanishes, i.e. $U=0$.
\end{lemma}
\begin{proof}
If $U\ne0$, it has a minimal support prime $\mathfrak q$.
Lemma~\ref{lem:tangential-finiteness} makes the relevant lengths finite,
and Lemmas~\ref{lem:page-inequality} and
\ref{lem:strict-support-inequality} contradict one another.
Thus $U = M/{x_1}M=0$.
\end{proof}
\AIadd{\formalref{\leandecl{Stafford38/Characteristic/CanonicalKoszulContradiction.lean}{83}{Stafford38.Characteristic.CanonicalKoszulContradiction.coordinate_cokernel_subsingleton},
\leandecl{Stafford38/Characteristic/CanonicalKoszulContradiction.lean}{47}{Stafford38.Characteristic.CanonicalKoszulContradiction.no_minimal_coordinate_cokernel_support}}}


\section{Characteristics of the characteristic variety}\label{sec:char}
\begin{corollary}\label{cor:charP}
Equip $Q$ with the order filtration $Q_m:=qF_m=\operatorname{im}(F_m\to Q)$
induced by the cyclic generator $q$ denoting the class of 1 in $Q$. Then
\begin{equation}\label{eq:char-contained}
  \Char(Q)\subseteq V(P)\subseteq T^*\bbA^n .
\end{equation}
\end{corollary}

\begin{proof}
$Q_\bullet$ is a good filtration because $Q$ is cyclic. Since $q\,\tilde d=[\tilde d]=0$ in
$Q$, the symbol $P=\sigma_N(\tilde d)$ annihilates the class $\bar q\in Q_0/Q_{-1}$
of $q$ in $\gr Q$: the element $q\tilde d\in Q_N$ is zero, and its image in
$Q_N/Q_{N-1}$ is $\bar q\cdot P$. Since $\gr Q$ is generated over $k_0[x,\xi]$
by $\bar q$, we get $P\in\Ann_{k[x,\xi]}\gr Q$, and
$\Char(Q)=V(\Ann\gr Q)\subseteq V(P)$ because $V$ reverses inclusions.
\end{proof}
\AIadd{\formalref{\leandecl{Stafford38/Characteristic/InitialIdeal.lean}{277}{Stafford38.CharacteristicInitialIdeal.canonical_orderPrincipalComponent_mem_initialIdeal},
\leandecl{Stafford38/Characteristic/InitialIdeal.lean}{287}{Stafford38.CharacteristicInitialIdeal.canonical_orderCharacteristicSupport_subset_principal_zeroLocus}}}
\begin{proposition}\label{prop:char-avoids}
The canonical quotient satisfies
\begin{equation}
 \Char(Q)\cap\pi^{-1}(H)=\varnothing,
 \label{eq:char-avoids}
\end{equation}
where $\pi:T^*\mathbb A^n\to\mathbb A^n$ is the cotangent projection\AIadd{;
equivalently, $\Ann_C\gr Q+(x_1)=C$, i.e.\ no prime of $\Supp_C\gr Q$ contains $x_1$}.
\end{proposition}
\begin{proof}
By Lemma \ref{lem:Uzero} $U = M/{x_1}M=0$. Since $M$ is finite over $C$, its quotient by ${x_1}$ has
support $\operatorname{Supp}_C(M)\cap V({x_1})$. Lastly, $\operatorname{Supp}_C(M) = \Char(Q)$, and $V({x_1})$ is the set of all points $\in k^{2n}$ where ${x_1} = 0$, which is nothing else than the pre-image of $H = \{{x_1} = 0\}$ under the projection, i.e. $V({x_1}) = \pi^{-1}(H)$. Since $M/{x_1}M = 0$, it has no support, so $\operatorname{Supp}(M/{x_1}M) = \operatorname{Supp}_C(M)\cap V({x_1}) =\Char(Q)\cap\pi^{-1}(H)= \varnothing$. 

%Formal adjoint changes $(x,\xi)$ to $(x,-\xi)$ and fixes $x$, so the same avoidance holds for the corresponding left module.
\AIadd{Here points are prime ideals of $C$: the argument is
$\Supp_C(M/x_1M)=\Supp_C M\cap V(x_1)$ for the finite $C$-module $M$.}
\end{proof}
\AIadd{\formalref[{stated for the support transported along $\xi\mapsto-\xi$, which fixes $x_1$; transported back verbatim}]{%
\leandecl{Stafford38/Characteristic/CanonicalKoszulContradiction.lean}{120}{Stafford38.Characteristic.CanonicalKoszulContradiction.canonical_support_avoidance},
\leandecl{Stafford38/Characteristic/SpecializedNoncharacteristicEquality.lean}{154}{Stafford38.SpecializedNoncharacteristicEquality.transposedSupport_disjoint_axis_iff}}}
\AIcomment{COMP-08}{\textbf{Information --- proof comparison; no edit requested.} The formal conclusion is scheme-theoretic: $\Supp(\gr Q)\cap V(x_1)=\varnothing$, equivalently $\Ann(\gr Q)+(x_1)=(1)$. It therefore persists under field extension. The formal statement uses the support transported by $\xi\mapsto-\xi$; this fixes $x_1$, so the avoidance statements are equivalent. \href{https://itpplasma.github.io/stafford38-formal/lean_proof_details.pdf\#nameddest=comparison-COMP-08}{Mathematical proof (COMP-08, p.~5).}}




\section{A conormal direction at infinity}\label{sec:conormal-review}
% add from here we work over C that means we define Q_c = ... etc
\AIadd{Throughout this section and the next, $\mathbb C$ may be replaced by any
algebraically closed field $K$ of characteristic zero; the formal statements are
proved in that generality, and Section~\ref{sec:proof} uses $K=\overline k$.}
Let $X=\mathbb A^n_{\mathbb C}$ be the affine $n$-space over $\mathbb{C}$. Because $\mathbb A^n_{\mathbb C}$ \AIreplace{is flat}{has a global coordinate frame}, $T^*X=X\times(\mathbb C^n)^*$ holds and we can therefore identify all cotangent fibres with $(\mathbb C^n)^*$. We write $[Z_0:\cdots:Z_n]$ for
homogeneous coordinates on the projective space $\mathbb P^n$, and embed
$X\hookrightarrow\mathbb P^n$ as the affine chart $\{Z_0\ne0\}$ via
$[1:x_1:\cdots:x_n]$, so that $x_i=Z_i/Z_0$ and the hyperplane at infinity
is $\{Z_0=0\}$ Let $Y\subset X$ be an irreducible closed subvariety of $X$. We define $\operatorname{Dir}(Y)$ to be the Zariski closure in  $\mathbb P((\mathbb C^n)^*)$  of the nonzero conormal vectors at all smooth affine points of $Y$ projected onto $\mathbb P((\mathbb C^n)^*)$. 
\AIadd{\formalref[{$Y$ given by a prime ideal; conormal space = annihilator of the Zariski tangent space; closure by homogeneous polynomials}]{%
\leandecl{Stafford38/Geometry/ProjectiveConormalDirections.lean}{37}{Stafford38.Geometry.ProjectiveConormalDirections.smoothConormalDirectionSet},
\leandecl{Stafford38/Geometry/ProjectiveConormalDirections.lean}{29}{Stafford38.Geometry.ProjectiveConormalDirections.projectiveHomogeneousClosure}}}


\begin{lemma}\label{lem:tangent-limit}
Suppose a formal arc in the projective closure of $Y$ has generic point in the
smooth affine locus, and put $m=\dim Y$. %needed so that tangent space is well defined
Let $L\subset\mathbb C[[t]]^{n+1}$ be a direct summand of rank $m+1$ whose generic
fibre is the affine cone over the embedded projective tangent $m$-plane. If
the reduction $L_0$ satisfies $L_0\subset\{Z_1=0\}$, then
$[dx_1]\in\operatorname{Dir}(Y)$.
\end{lemma}

\begin{proof}
To prove this lemma, we will construct an arc such that its generic fibre is a conormal direction at a smooth point of $Y$ (hence lies in $\operatorname{Dir}(Y)$), identify its special fibre as $[dx_1]$ and conclude from the closedness of $\operatorname{Dir}(Y)$ that it must also lie in $\operatorname{Dir}(Y)$. 
The annihilator $L^\perp$ of $L$ is also a direct summand, and therefore reduces onto $L_0^\perp$. 
The dual basis vector $Z_1^\vee$ annihilates $L_0$ since $L_0\subset \{Z_1 = 0\} = \ker(Z_1^\vee)$, and hence $Z_1^\vee \in L_0^\perp$. Because the reduction is a surjective map, we can lift $Z_1^\vee\in L_0^\perp$ to $\ell(t)\in L^\perp$, satisfying $\ell(0) = Z_1^\vee.$ 
At the generic point, $\ell(t)$ annihilates the generic fibre of $L$ and hence by hypothesis the cone over the embedded projective tangent space. The fibre coordinates $(\ell_1(t),\ldots,\ell_n(t))$ of $\ell(t)$ in \AIreplace{the the}{the}
affine chart $Z_0\ne0$ therefore form a nonzero affine conormal
covector. Because $\ell(0)=Z_1^\vee\neq 0$, the reduction of the fibre coordinates at $t = 0$ is nonzero, so their projective direction extends over $\mathbb C[[t]]$. Specifically, this reduction is $(1,0,\ldots,0)$, giving the special fibre $[dx_1]$. %meaning evaluating the projection at t = 0 gives you [dx_1].
\AIreplace{Since $\operatorname{Dir}(Y)\subset \mathbb P((\mathbb C^n)^*)$ is closed, hence proper over $\mathbb C$, $[dx_1]$ must lie in $\operatorname{Dir}(Y)$ by the valuative criterion \mbox{\cite[Theorem II.4.7]{hartshorne}}.}{The generic direction lies in $\operatorname{Dir}(Y)$: every homogeneous polynomial vanishing on $\operatorname{Dir}(Y)$ vanishes on the conormal directions of the $\mathbb C((t))$-point $\gamma(\eta)$ as well (the conormal variety over $Y_{\mathrm{sm}}$ is defined over $\mathbb C$, and its $\mathbb C$-points are dense in it). Hence the defining equations of the closed set $\operatorname{Dir}(Y)$ vanish on the direction of $\ell(t)$ in $\mathbb C((t))$, hence on its value at $t=0$, which is $[dx_1]$. No valuative criterion is needed.}
\end{proof}
\AIcomment{GEO-04}{\textbf{Proposed change --- author acceptance pending.} The printed proof did not justify that a conormal direction
at the $\mathbb C((t))$-point $\gamma(\eta)$ belongs to $\operatorname{Dir}(Y)$,
which is defined from $\mathbb C$-points; the red text supplies the
scalar-extension step, which the formalization proves as
\texttt{scalarExtension\_vanishing}.}
\AIadd{\formalref[{any algebraically closed field and any axis; auxiliary, not used by the formal proof of Theorem~\ref{thm:asymptotic-conormal}}]{%
\leandecl{Stafford38/Geometry/GeneralTangentLimitCriterion.lean}{497}{Stafford38.Geometry.GeneralTangentLimitCriterion.tangent_limit_criterion_of_directSummand};
scalar extension \leandecl{Stafford38/Geometry/ConormalScalarExtensionVanishing.lean}{134}{Stafford38.Geometry.ConormalScalarExtensionVanishing.scalarExtension_vanishing}}}
\begin{AIaddition}{The generic conormal direction and the passage to $t=0$}
Let $\eta$ be the generic point of the arc, and write
$\gamma(\eta)=[v_0:\cdots:v_n]$ with $v_i\in\mathbb C((t))$ and $v_0\neq0$, since the generic point is affine. Put
$y=(v_i/v_0)_{i\ge1}\in Y(\mathbb C((t)))$ and let $T_yY$ be the Zariski
tangent space of $I(Y)\otimes\mathbb C((t))$ at $y$.

\emph{(1) Dehomogenization.} For $\tau\in T_yY$, the vector $(0,v_0\tau)$ lies in the
affine cone over the projective tangent plane, which is the generic fibre of $L$.
Since $\ell(t)$ annihilates that fibre,
$0=v_0\sum_{i\ge1}\ell_i\tau_i$. Hence $\ell'=(\ell_1,\dots,\ell_n)$ annihilates
$T_yY$, and so $\ell'\in(T_yY)^\perp=\operatorname{span}\{\nabla f(y):f\in I(Y)\}$,
i.e.\ $(y,\ell'(t))$ is a $\mathbb C((t))$-point of the equation-conormal locus
$\mathcal C_{\mathbb C((t))}$ obtained by extending the equation-conormal locus of $Y$ to $\mathbb C((t))$. Smoothness of $y$ is not used.

\emph{(2) Passage to $\mathbb C$-points.} $\Dir(Y)$ is defined through conormal
vectors at $\mathbb C$-points, and $y$ is not one. Let $P$ be a homogeneous
polynomial vanishing at every direction in $\Dir(Y)$. Then $\xi_1P$ vanishes on
$N(Y)=\bigcup_{z\in Y_{\sm}}N^*_zY$. By generic smoothness, this vanishing
extends to the equation-conormal locus: a polynomial in gradient covectors
vanishing on the dense smooth locus vanishes on all of $Y$. Parametrize
covectors by finite linear combinations of gradients and apply the
Nullstellensatz in $\mathbb C[x,\lambda]$; the vanishing then persists
under extension to $\mathbb C((t))$. Thus
$\ell_1(t)P(\ell'(t))=0$. Since $\ell_1(0)=1$, $\ell_1$ is a unit of
$\mathbb C[[t]]$, hence $P(\ell'(t))=0$ in $\mathbb C[[t]]$.

\emph{(3) Closedness replaces the valuative criterion.} The coefficients of $P$
are constants and $\ell'(t)\in\mathbb C[[t]]^n$, so the constant term of
$P(\ell'(t))$ is $P(\ell'(0))=P(1,0,\dots,0)$, and it vanishes. As $P$ was an
arbitrary homogeneous equation of the closed set $\Dir(Y)$, we get
$[dx_1]\in\Dir(Y)$. No properness is needed, because the limit point
$\ell'(0)=e_1$ is given explicitly.
\end{AIaddition}
\AIadd{\formalref[{generic point only Zariski-tangent; no smoothness needed for this step}]{%
\leandecl{Stafford38/Geometry/ConormalScalarExtensionVanishing.lean}{134}{Stafford38.Geometry.ConormalScalarExtensionVanishing.scalarExtension_vanishing};
\leandecl{Stafford38/Geometry/SmoothConormalFibreVanishing.lean}{21}{Stafford38.Geometry.SmoothConormalFibreVanishing.fibreLift_mem_vanishingIdeal_equationConormal};
\leandecl{Stafford38/Geometry/LaurentConormalDirection.lean}{68}{Stafford38.Geometry.LaurentConormalDirection.residue_mem_extensionFibreClosure_of_laurent_generic};
\leandecl{Stafford38/Geometry/GeneralTangentLimitCriterion.lean}{368}{Stafford38.Geometry.GeneralTangentLimitCriterion.exists_axis_laurent_smooth_conormal_direction}}}
\AIcomment{COMP-10}{\textbf{Information --- proof comparison; no edit requested.} The formal tangent-limit criterion is auxiliary. It assumes a generically smooth arc and a direct-summand lattice whose generic span is the lifted tangent space, with the specified axis component vanishing on its residue. Its proof does not use the smoothness hypothesis or the stated rank, but these remain assumptions of the formal statement. The formal proof of Theorem~\ref{thm:asymptotic-conormal} follows a separate Laurent-series route, so it does not certify every step of the printed tangent-limit argument. \href{https://itpplasma.github.io/stafford38-formal/lean_proof_details.pdf\#nameddest=comparison-COMP-10}{Mathematical proof (COMP-10, p.~5).} See ROUTE-01 for the separate author decision on the printed geometric proof.}

\begin{theorem}\label{thm:asymptotic-conormal}
If $Y$ is nonempty and $Y\cap H=\varnothing$, then
\[[dx_1]\in\operatorname{Dir}(Y).\]
\end{theorem}
\AIadd{\formalref[{over every algebraically closed field of characteristic zero; the formal proof follows a different route, written out after the printed proof}]{%
\leandecl{Stafford38/Geometry/GeneralAsymptoticConormal.lean}{63}{Stafford38.Geometry.GeneralAsymptoticConormal.coordinate_axis_mem_projective_conormal_directions},
affine-cone form \leandecl{Stafford38/Geometry/GeneralAsymptoticConormal.lean}{28}{Stafford38.Geometry.GeneralAsymptoticConormal.coordinate_axis_mem_smooth_fibre_closure}}}

\AIcomment{ROUTE-01}{\textbf{Action required --- approve the publication proof.} The proposed publication route is the machine-verified Laurent-series argument in comparison COMP-11. Max should check that mathematical account and its statement, not re-audit the tangent-lattice alternative below. The printed alternative is retained for the authors; it enters the publication review only if they decide to publish it and provide its own formal correspondence. The theorem's machine verification does not certify that alternative's intermediate steps.}
\begin{proof}
If $Y$ is zero-dimensional, its tangent space is trivial, so every cotangent direction occurs in its conormal
fibre. If $\dim Y>0$, but the coordinate function $x_1|_Y$ is constant on $Y$, it is a nonzero constant (since $Y\cap \{x_1=0\} = \varnothing$) and $dx_1$ annihilates every tangent space and is therefore a conormal vector to $Y$, i.e. $[dx_1]\in \operatorname{Dir}(Y)$. We may therefore assume that $m=\dim Y>0$ and that $x_1|_Y$ is nonconstant.

Let $\overline Y\subset\mathbb P^n$ be the closure of $Y$ in $\mathbb P^n$ and let
$\nu:\widetilde Y\to\overline Y$ be its normalization. Normalization is
finite in this setting \cite[\AIreplace{Tag 0BXQ}{Tag~0BXR}, Lemma~33.27.1]{Stacks}. 
The rational function $x_1=Z_1/Z_0$ has a zero prime divisor $D$\AIadd{ }on
$\widetilde Y$, , i.e. a prime divisor with $\ord_D(x_1)>0$. Otherwise $x_1^{-1}$ would have no pole at any codimension-one point of $\widetilde Y$. Since $\widetilde Y$ is normal,
$x_1^{-1}$ would be globally regular \cite[Tag 031T]{Stacks}. 
As $\widetilde Y$ is finite (hence proper) over $\overline Y$, it is proper over $\mathbb C$.
\AIadd{Finite normalization and preservation of projectivity are standard\cite[\AIreplace{Tag~0BXQ}{Tags~0BXR, 0GK4}, Lemmas~33.27.1--2]{Stacks}.} By \cite[\AIremove{Tag 04L0}\AIadd{Tag 04L2}, Lemma~33.26.2]{Stacks}, any globally regular function must be constant. This contradicts our standing assumption that $x_1|_Y$ is nonconstant.
Since $Y\cap \{x_1=0\}=\varnothing$, $x_1$ \AIreplace{does not vanish on $Y$}{is a unit on $Y$ (Nullstellensatz)}, so the generic point
of $D$ lies outside $\nu^{-1}(Y)$, hence in the hyperplane $\{Z_0 = 0\}$ at infinity. %Because we defined $Y$ to be closed in $X$, $\overline{Y}\backslash Y$ cannot lie in $X$, hence at infinity, $\nu^{-1}(\overline{Y}\backslash Y)$. %above the hyperplane 
%The rational function \[ x_1=Z_1/Z_0\] has a zero prime divisor $D$ on $\widetilde Y$, i.e. an irreducible closed subvariety of codimension 1 such that $x_1$ vanishes on $D$. Otherwise $x_1^{-1}$ would have no pole at any codimension-one point. Normality would make it globally regular \cite[Tag 031T]{Stacks}, and projectivity would make it constant \cite[Tag 04L0, Lemma~33.26.2]{Stacks}. 
We choose $Z_j$ nonzero at the generic point of $D$, and define
$q_i=Z_i/Z_j$. With respect to the divisorial valuation at $D$, we define
\begin{equation}a=\operatorname{ord}_D(q_0),\qquad
 b=\operatorname{ord}_D(q_1).
\label{eq:ab}\end{equation}
Indeed, because the generic point lies in $\{Z_0=0\}$, $a$ must be positive\AIadd{; as $a,b>0$ while $q_j=1$, the index $j$ differs from $0$ and $1$}. $b$ is finite since $Z_1$
does not vanish identically on $\overline Y$, and we recall $x_1=Z_1/Z_0=q_1/q_0$, therefore $\ord_D(x_1) = \ord_D(q_1)-\ord_D(q_0) = b-a$. Because $D$ was specifically chosen to be a zero prime divisor of $x_1$, $\ord_D(x_1)>0$, hence $b>a.$
We shrink a neighbourhood $U$ of the generic point of $D$ so that $Z_j$ is
invertible. Defining $D\cap U=(s)$, we can find regular functions $w_0$ and $w_1$ that do not vanish identically on $D$, such that
\[ q_0=s^aw_0,\qquad q_1=s^bw_1,\]
Restricting the normalization to $D$ inherits finiteness, so at the generic point of $D$, the map has finite fibers. 
Because $\nu_{|D}$ is finite, its function-field extension is separable since $\mathbb C$ is a field of characteristic 0 (\cite[Tag 0322]{Stacks}), and consequently, $\nu_{|D}$ is smooth at the generic point of $D$ \cite[Tag 07ND]{Stacks}\AIadd{, applied to the dominant finite morphism $D\to\nu(D)$ onto its image}. 
 Since smoothness is an open condition, $\nu_{|D}$
is smooth on a dense open subset of $D$. Equivalently, since $D$ and $\nu(D)$ have equal dimension
$m-1$, its differential has rank $m-1$.
\AIcomment{GEO-01}{\textbf{Proposed change --- author acceptance pending.} The generic-smoothness step is valid for the dominant
finite map $D\to\nu(D)$ (Tag~07ND). Name that image as the target and
choose the point also over its smooth locus; then the differential into
the projective chart has the asserted rank $m-1$.}

Choose a closed point $p$ of $D\cap U$ in the smooth loci of both $D$ and
$\widetilde Y$, where the differential of $\nu_{|D}$ has rank $m-1$, and
outside $V(w_0w_1)$, i.e. $w_0\neq 0$ and $w_1\neq 0$ at the chosen point\AIadd{, and lying over the smooth locus of $\nu(D)$}. Normality of $\widetilde Y$ ensures that its singular locus does not contain $D$ and because all conditions hold on a dense open subset of $D$, their finite intersection is nonempty, so such a point exists. 

%tag 0C0S Lemma 10.160.10 says that completed local ring is isomorphic to a power series ring over its residue field
We can define a completed local ring at $p$, which is isomorphic to a power series ring \cite[Tag 0C0S, Lemma 10.160.10]{Stacks}. We can choose a formal coordinate system centered around $p$ with parameters
$(t,z_2,\ldots,z_m)$. We choose $t=s$: since our point lies in the smooth locus of $D = \{s = 0\}$, $ds\neq 0$ there, so $s$ extends to a regular system of parameters. 

The set $\widetilde Y\setminus\nu^{-1}(Y_{\mathrm{sm}})$ (with $Y_{\mathrm{sm}}$ being the smooth locus of $Y$) is a proper closed subset of $\widetilde Y$ containing $p$. Let $I\subset
\widehat{\mathcal O}_{\widetilde Y,p}\cong\mathbb C[[t,z_2,\ldots,z_m]]$ be
the ideal of this closed set in the completed local ring, so that every element of $I$ vanishes at every point of this closed set, and choose any
nonzero $f\in I$. We expand $f$ in homogeneous parts,
\[ f = f_d(t,z_2,\ldots,z_m) + f_{d+1}(t,z_2,\ldots,z_m) + \cdots,\]
where $f_d\ne0$ is the lowest-degree nonzero homogeneous part. Since
$f_d(1,z_2,\ldots,z_m)$ is a nonzero polynomial in $z_2,\ldots,z_m$, it
vanishes only on a proper closed subset of $\mathbb C^{m-1}$. We now choose
constants $c_2,\ldots,c_m$ outside this subset, so that
\[f_d(1,c_2,\ldots,c_m)\ne0.\] When we change the coordinate system by replacing $z_\alpha$ by
$z_\alpha-c_\alpha t$, we therefore guarantee that $f\ne0$ for $t\ne0$
along the (tilted) formal $t$-axis. Since $f$ vanishes on all of
$\widetilde Y\setminus\nu^{-1}(Y_{\mathrm{sm}})$ (as $f$ lies in its
ideal), this forces every point of the arc with $t\ne0$ -- in particular
its generic point -- to lie outside $\widetilde Y\setminus
\nu^{-1}(Y_{\mathrm{sm}})$, i.e.\ in $\nu^{-1}(Y_{\mathrm{sm}})$.

Setting $z_2=\cdots=z_m=0$ in these coordinates therefore defines a
transverse map $\gamma:\operatorname{Spec}\mathbb C[[t]]\to\widetilde Y$
whose generic point lies, by construction, in $Y_{\mathrm{sm}}$, the smooth
affine locus of $Y$.

In the new coordinates of the local ring, we can accordingly rewrite
\begin{equation} q_0=t^au_0,\qquad q_1=t^bu_1,
 \label{eq:q-orders}\end{equation}
where $u_0,u_1$ are units since $w_0(p)\neq 0, w_1(p)\neq 0$. 

In the coordinates $(t,z_2,\ldots,z_m)$, each $q_i$ is a regular function
on $\widetilde Y$ near $p$, so $v=(q_0,\ldots,q_n)^\top$, with $q_j = 1$  is a local
parametrization, in these coordinates, of $\nu$'s image in the chart
$Z_j\ne0$.
We define the partial derivative of $v$ with respect to $t$ and the $z_\alpha$
\[ V_\alpha=\frac{\partial v}{\partial z_\alpha},\qquad
 T=\frac{\partial v}{\partial t}. \]
$V_\alpha$ vanish in rows $0$, $1$, and $j$ at $t = 0$, since $q_j = 1$ is constant, and $q_0$ and $q_1$ (and hence their derivatives) have factors $t$ vanishing at $t=0$. Since locally, $D = \{t=0\}$, they are
the images of a tangent basis of $D$ under the differential of $\nu_{|D}.$ Because we chose $p$ such that $\nu_{|D}$ would have full rank and is hence injective, the images $V_\alpha$ of the tangent basis of $D$ are also independent. 
Since the reductions of $V_2,\ldots,V_m$ are independent and vanish in
rows $0,1,j$, they have rank $m-1$, so there exist $m-1$ further
coordinate rows $i_1,\ldots,i_{m-1}$ (necessarily distinct from $0,1,j$) in
which the corresponding $(m-1)\times(m-1)$ submatrix of $(V_2,\ldots,V_m)$
is invertible. Because this minor is invertible at $t=0$, it remains
invertible over $\mathbb C[[t]]$, so there are unique coefficients
$\lambda_2(t),\ldots,\lambda_m(t)\in\mathbb C[[t]]$ solving
\[\sum_{\alpha=2}^m \lambda_\alpha(t)\,(V_\alpha)_{i_k} = T_{i_k},
 \qquad k=1,\ldots,m-1.\]
Setting
\[ T' := T - \sum_{\alpha=2}^m \lambda_\alpha(t)\,V_\alpha\]
therefore gives a column vector whose entries in rows
$i_1,\ldots,i_{m-1}$ vanish identically by construction. The entries of $V_\alpha$ in rows $0$ and $1$ have valuation at least $a$ and $b$ respectively, whereas
\[ \partial_tq_0=a t^{a-1}u_0+O(t^a),\qquad
 \partial_tq_1=b t^{b-1}u_1+O(t^b).\]
Because we work over a field of characteristic zero, the displayed
leading coefficients are nonzero, so
\begin{equation} \operatorname{ord}_t(T_0)=a-1,\qquad \operatorname{ord}_t(T_1)=b-1.
 \label{eq:T-orders-raw}\end{equation}
Since the entries of $V_\alpha$ in rows $0,1$ have strictly higher
valuation, subtracting a combination of $V_2,\ldots,V_m$ from $T$ to form
$T'$ cannot change these orders, hence
\begin{equation} \operatorname{ord}_t(T'_0)=a-1,\qquad\operatorname{ord}_t(T'_1)=b-1. 
\label{eq:T-orders}\end{equation}

Let $c$ be the minimum valuation of the entries of $T'$. Then $c\leq a-1$,
and $\widehat T=t^{-c}T'$ is regular and its evaluation at $t=0$ is nonzero. Moreover
\begin{equation}\operatorname{ord}_t(\widehat T_1)\geq b-a>0.
 \label{eq:T1-vanish}\end{equation}
 
Since $\gamma$'s generic point lies in $Y_{\mathrm{sm}}$, and
normalization is an isomorphism over the smooth locus,
\AIadd{by the finite birational normal-target criterion\cite[Tag~0AB1]{Stacks},} $\widetilde Y$ and
$Y$ agree there; working over $\mathbb C((t))$, we may therefore regard
$v$ as a local parametrization of $Y$ itself near $\gamma$'s
generic point.
Since $(t,z_2,\ldots,z_m)$ is a full coordinate system on the
$m$-dimensional $\widetilde Y$, the $m$ vectors $T,V_2,\ldots,V_m$ form a
tangent basis at that point; together with $v$ itself, they span the
affine cone over the projective tangent $m$-plane to $Y$ along $\gamma$.
Since $T'$ differs from $T$ only by a combination of $V_2,\ldots,V_m$, and
$\widehat T$ differs from $T'$ only by multiplication by the invertible
scalar $t^{-c}\in\mathbb C((t))^\times$, replacing $T$ by $T'$ and then by
$\widehat T$ does not change this span.

Also evaluated at $t = 0$, the columns $v,V_2,\ldots,V_m,\widehat T$ are independent vectors. Indeed,
$v_j=1$, while the other vectors evaluated at $t=0$ have $j$th coordinate zero. $V_\alpha(0)$ were chosen to be independent among themselves, and by construction $\widehat T(0)\neq 0$. \AIreplace{If $\widehat T(0)$ were a combination of $v(0),V_2(0),\ldots,V_m(0)$, matching the zero entries in rows $i_1,\ldots,i_{m-1}$ would
force the $V_\alpha(0)$-coefficients to vanish (by invertibility of the
minor), and matching row $j$ would then force the $v(0)$-coefficient to
vanish too (since $v_j(0)=1\ne0=\widehat T_j(0)$), contradicting that
$\widehat T(0)$ is nonzero, hence also $\widehat T(0)$ is independent from the other vectors.}{If $\widehat T(0)$ were a combination of $v(0),V_2(0),\ldots,V_m(0)$, matching row $j$ first forces the $v(0)$-coefficient to vanish (since $v_j(0)=1$ while all other vectors have $j$th entry $0$); then matching the rows $i_1,\ldots,i_{m-1}$, where $\widehat T(0)$ vanishes, forces the $V_\alpha(0)$-coefficients to vanish by invertibility of the minor. So $\widehat T(0)=0$, a contradiction.}
$v, \widehat T, V_2,\dots, V_m$ therefore span the direct-summand tangent
lattice $L$ required by Lemma~\ref{lem:tangent-limit} and their evaluations at $t = 0$ span the reduction $L_0$. 
As a remark, we note that $L_0$ is the Grassmannian limit.
\AIcomment{GEO-02}{\textbf{Proposed change --- author acceptance pending.} In the independence argument above, use row $j$
first: it forces the coefficient of $v(0)$ to be zero. Only then do the
selected minor rows force all $V_\alpha(0)$ coefficients to vanish.
Otherwise the selected rows may contain a contribution from $v(0)$.
Nonzero $\widehat T(0)$ gives the contradiction and a unit maximal minor
gives the direct-summand lattice.}
All vectors $v, \widehat T, V_2,\dots, V_m$ evaluated at $t=0$
have zero $Z_1$-coordinate, so
\[ L_0\subset\{Z_1=0\}.\]
By Lemma~\ref{lem:tangent-limit} $[dx_1] \in \operatorname{Dir}(Y).$
\end{proof}

\AIcomment{COMP-11}{\textbf{Information --- proof comparison; no edit requested.} The formal proof works over every algebraically closed field of characteristic zero. It replaces the tangent-limit route by a Laurent-series argument. The constant-coordinate case is treated separately; otherwise a divisorial valuation gives $dx_1=\sum_j c_j\,dx_j$ with each $c_j$ of positive valuation. The conormal covector $e_1-\sum_jc_je_j$ therefore tends to $e_1$, and the equation-conormal closure and scalar-extension arguments give the required direction. The detailed valuation and residue argument is in the technical supplement; it does not justify a gap in a different printed proof. \href{https://itpplasma.github.io/stafford38-formal/lean_proof_details.pdf\#nameddest=comparison-COMP-11}{Mathematical proof (COMP-11, p.~6).} See ROUTE-01 for the separate author decision on the printed geometric proof.}
\AIadd{\formalref[{any algebraically closed field of characteristic $0$; (i) is the affine-cone form used in Section~\ref{sec:involutive-review}}]{%
\leandecl{Stafford38/Geometry/GeneralAsymptoticConormal.lean}{63}{Stafford38.Geometry.GeneralAsymptoticConormal.coordinate_axis_mem_projective_conormal_directions};
\leandecl{Stafford38/Geometry/GeneralDivisorialVisibleFrame.lean}{42}{Stafford38.Geometry.GeneralDivisorialVisibleFrame.generalDivisorialVisibleFrameExistence};
\leandecl{Stafford38/Geometry/DivisorTangentLattice.lean}{326}{Stafford38.Geometry.DivisorTangentLattice.VisibleDivisorFrame.exists_maximalIdeal_coefficients};
\leandecl{Stafford38/Geometry/ConormalScalarExtensionVanishing.lean}{134}{Stafford38.Geometry.ConormalScalarExtensionVanishing.scalarExtension_vanishing}}}

\section{The involutive obstruction}\label{sec:involutive-review}
\begin{definition}[Fibre-conical]\label{def:fibre-can}
A closed subset $W\subseteq T^*X$ is \emph{fibre-conical} if it is stable under the fibrewise scaling $(y,\xi)\mapsto(y,t\xi)$, $t\in\mathbb{G}_m$. 
\end{definition}

\begin{theorem}[Involutive exclusion]\label{thm:coisotropic-exclusion} %coistropic is equivalent with involutive but involutive is what is actually used
Let $W\subset T^*X$ be a reduced closed fibre-conical and involutive subset. Let
$P\in \mathbb C[\xi_1,\ldots,\xi_n]$ be homogeneous and suppose
$P(1,0,\ldots,0)\neq0$.
There is no nonempty $W$ satisfying
\[
 W\cap\pi^{-1}(H)=\varnothing,\qquad W\subset V(P).
\]
\end{theorem}
\begin{AIaddition}{Formal statement of Theorem~\ref{thm:coisotropic-exclusion}}
Lean proves the following strengthening. Let $K$ be an algebraically closed field of characteristic zero and $n\ge1$.
Let $W\subset K^n\times K^n$ be nonempty and Zariski closed, with
$(y,a\xi)\in W$ whenever $(y,\xi)\in W$ and $a\in K^\times$, and with
$\{I(W),I(W)\}\subseteq I(W)$. Let $P\in K[\xi_1,\dots,\xi_n]$ be \emph{any}
polynomial, not necessarily homogeneous, with $P(\xi)=0$ for all $(y,\xi)\in W$ and
$P(1,0,\dots,0)\neq0$. Then some $(y,\xi)\in W$ has $y_1=0$, i.e.\
$W\cap\pi^{-1}(H)\neq\varnothing$. The proof uses only the base-relative
fragment $\{f\circ\pi,I(W)\}\subseteq I(W)$ for $f\circ\pi\in I(W)$. The Lean bracket is
the negative of the one in the paper, which does not affect involutivity.
``Involutive'' means $\{I(W),I(W)\}\subseteq I(W)$, not that $I(W)$ is a Poisson ideal.
For $K=\mathbb C$ this contains Theorem~\ref{thm:coisotropic-exclusion}.
\end{AIaddition}
\AIadd{\formalref[{any algebraically closed field of characteristic $0$; $P$ need not be homogeneous}]{%
\leandecl{Stafford38/Geometry/GeneralCoisotropicSets.lean}{108}{Stafford38.Geometry.GeneralCoisotropicSets.exists_zero_base_coordinate_of_isFibreConical};
\leandecl{Stafford38/Geometry/GeneralCoisotropicExclusion.lean}{29}{Stafford38.Geometry.GeneralCoisotropicExclusion.exists_zero_base_coordinate};
\leandecl{Stafford38/Characteristic/BaseRelativePoisson.lean}{28}{Stafford38.Characteristic.BaseRelativePoisson.IsBaseRelativePoisson};
\leandecl{Stafford38/Geometry/GeneralCoisotropicSetsTest.lean}{100}{Stafford38.Geometry.GeneralCoisotropicSetsTest.exact_complex_manuscript_coisotropic_consumer}}}

\begin{proof}
Supposing that such a $W$ exists, we define $Y=\overline{\pi(W)}$. We write $I(Y)$
and $I(W)$ for the radical vanishing ideals, and define
$W_y=W\cap T_y^*X$ (using the identification $T_y^*X\cong (\mathbb C^n)^*$ as before). Closedness and
fibre-conicality imply
\begin{equation}
 Y\times\{0\}\subset W.
 \label{eq:zero-section}
\end{equation}
Every nonempty conical fibre contains its origin, so in fact
$\pi(W)=\{y:(y,0)\in W\}$ is closed and equals $Y$. In particular,
$Y\cap H=\varnothing$.

For $f\in I(Y)$, its pullback belongs to $I(W)$. Involutivity gives
$\{f,I(W)\}\subset I(W)$. The Hamiltonian derivation of $f$ is
\[
 H_f=\{-,f\}=\sum_i\frac{\partial f}{\partial x_i}
 \frac{\partial}{\partial\xi_i}.
\]
Due to the involutivity, $H_f$ preserves $W$.
Since $H_f$ involves only $\partial/\partial\xi_i$, it strictly lowers
$\xi$-degree, hence is locally nilpotent on $\mathbb C[x,\xi]$; its
exponential is therefore a well-defined, finite sum on any polynomial,
equal to vertical translation by $t\,df$: for any point $(y, \xi) \in T^*X$, solving the flow equations
$\dot y=0$, $\dot\xi=\partial_yf$ gives $(y,\xi)\mapsto(y,\xi+t\,df_y)$.
These translations commute, so starting from $(y,0)\in W$ (eq. \eqref{eq:zero-section}) and composing them for varying $f\in I(Y)$ and $t\in\mathbb C$ shows that every point $(y,\sum_it_i\,d(f_i)_y)$ lies in
$W$, i.e. $\operatorname{span}\{df_y: f \in I(Y)\}\subset W_y$. %technicality: eigentlich W_y = {(y, xi) with fixed y such that (y, xi) in W} but since y is constant there is a "canonical identification T_y^*X \cong C^n* so we can just drop the y component.

Let $Y_0$ be an irreducible component of $Y$. At a general smooth point $y$ of
$Y_0$ outside the other components, $I(Y)_y=I(Y_0)_y$, so at this point, $\operatorname{span}\{df_y: f \in I(Y)\} = \operatorname{span}\{df_y: f \in I(Y_0)\} = N_y^*Y_0$, the conormal space to $Y_0$. Since $N_y^*Y_0\subset W_y$ holds at a dense open set of points in the smooth locus $y\in (Y_0)_{\mathrm{sm}}$ and $W$ is closed, the closure of the conormal bundle over the smooth locus $T^*_{Y_0}X := \overline{\bigcup_{y\in(Y_0)_{\mathrm{sm}}}N^*_yY_0}$ lies in $W$
\begin{equation} T^*_{Y_0}X\subset W.
 \label{eq:conormal-in-W}\end{equation}
Since $Y_0$ is an irreducible component of the closed set $Y$, $Y_0$ is nonempty, irreducible and itself closed. Furthermore, from $Y\cap H = \varnothing$ follows $Y_0\cap H = \varnothing$, so the \AIreplace{hypothese s}{hypotheses} of
Theorem \ref{thm:asymptotic-conormal} are fulfilled and hence $[dx_1] \in \operatorname{Dir}(Y_0).$ Since $\operatorname{Dir}(Y_0)$ is, by definition, the projectivization of $T^*_{Y_0}X$ (the closure of the
conormal directions over the smooth locus of $Y_0$), this places $[dx_1]$
in the closure of the projective fibre directions of $T^*_{Y_0}X$, hence
of $W$ by \eqref{eq:conormal-in-W}.
Since $P$ is homogeneous in the fibre variables $\xi$ alone, the
condition $P(\xi)=0$ depends only on the direction $[\xi]$, not on the
base point $x$ or the scale of $\xi$; it therefore defines a well-defined
closed hypersurface $V_+(P)\subset\mathbb P((\mathbb C^n)^*)$, the
projectivized vanishing locus of $P$. As $W\subset V(P)$, every
projectivized fibre direction of $W$ -- regardless of its base point, even
one escaping to infinity -- lies in this same $V_+(P)$. \AIreplace{Since $[dx_1]$ is a projectivized fibre direction of $W$, $[dx_1]\in V_+(P)$,}{Every conormal direction at a smooth point of $Y_0$ is a fibre direction of $W$ and hence lies in the closed set $V_+(P)$; therefore so does their closure $\operatorname{Dir}(Y_0)\ni[dx_1]$. Thus $[dx_1]\in V_+(P)$,} or equivalently $P(1,0,\ldots,0)=0$, a contradiction to the theorem's hypotheses. Therefore, no such $W$ that is nonempty can exist.
\AIcomment{GEO-03}{\textbf{Proposed change --- author acceptance pending.} The preceding result places $[dx_1]$ in the
\emph{closure} of projective conormal directions, not necessarily in a
fibre of $W$ at a finite base point. Every conormal direction lies in
$V_+(P)$; that set is closed, so the limiting direction does too.
This yields the same contradiction $P(1,0,\ldots,0)=0$.}
\end{proof}
\AIcomment{COMP-13}{\textbf{Information --- proof comparison; no edit requested.} The formal proof first excludes $x_1$ from every minimal base prime by iterated Poisson brackets with $P$. Assuming avoidance of $H$, it then forces those primes to be homogeneous in the transverse base coordinates and uses an asymptotic conormal direction to contradict $P(1,0,\ldots,0)\ne0$. This is an alternative route, not a verification of each intermediate step above. \href{https://itpplasma.github.io/stafford38-formal/lean_proof_details.pdf\#nameddest=comparison-COMP-13}{Mathematical proof (COMP-13, p.~8).} See ROUTE-01 for the separate author decision on the printed geometric proof.}
\AIadd{\formalref{%
\leandecl{Stafford38/Geometry/GeneralCoisotropicExclusion.lean}{29}{Stafford38.Geometry.GeneralCoisotropicExclusion.exists_zero_base_coordinate};
\leandecl{Stafford38/Characteristic/BaseRelativePoisson.lean}{71}{Stafford38.Characteristic.BaseRelativePoisson.zeroSection_stable_under_differential_translation_of_isBaseRelativePoisson};
\leandecl{Stafford38/Geometry/GeneralComponentConormalContainment.lean}{76}{Stafford38.Geometry.GeneralComponentConormalContainment.equationConormalClosure_minimalPrime_subset_zeroLocus};
\leandecl{Stafford38/Geometry/GeneralAsymptoticConormal.lean}{28}{Stafford38.Geometry.GeneralAsymptoticConormal.coordinate_axis_mem_smooth_fibre_closure}}}
\AIcomment{GEO-05}{\textbf{Optional refinement --- not required.} Optional statement refinements suggested by the formal
version, for the authors to decide: in Section~\ref{sec:conormal-review} say
that $Y$ is reduced and irreducible (so $I(Y)$ is prime) and identify
$(\mathbb C^n)^*$ with $\mathbb C^n$ via $dx_i$; in Lemma~\ref{lem:tangent-limit}
the arc is $\mathbb C[[t]]$-valued with affine generic point, and ``smooth'' can be
weakened to ``Zariski tangent space''; in Theorem~\ref{thm:asymptotic-conormal}
``$Y$ nonempty'' is automatic, and the affine-cone form (i) of the red proof is
what Theorem~\ref{thm:coisotropic-exclusion} actually uses; in
Theorem~\ref{thm:coisotropic-exclusion} ``reduced'' is vacuous for a point set,
the last sentence should read ``if $W\neq\varnothing$ and $W\subset V(P)$, then
$W\cap\pi^{-1}(H)\neq\varnothing$'', and homogeneity of $P$ can be dropped.}

\section{Proof of Theorem~\ref{thm:main}}\label{sec:proof}
\AIcomment{ALG-05}{\textbf{Proposed change --- author acceptance pending.} \emph{Updated 2026-09-30: the red text below now works over
$\overline k$ as the formalization does, states the base-change lemma and the
fibre-conicality, and replaces the placeholder citations.} For descent, flat field extension gives
$\gr(Q\otimes_{k_0}\mathbb C)\cong(\gr Q)\otimes_{k_0}\mathbb C$.
Empty support of the finite graded module implies it is zero;
exhaustiveness and $F_{-1}Q=0$ then imply $Q_{\mathbb C}=0$.
Faithful flatness detects $Q=0$\cite[Tag~00HP]{Stacks}. The placeholder
citations below are retained visibly as proposed removals.}

\AIadd{Let $K=\overline k$ and $Q_K:=Q\otimes_kK=A_n(K)/(\tilde dA_n(K)+x_1^N\tilde dA_n(K))$
(right exactness of $\otimes$). The order filtration of $Q_K$ is the base change of
that of $Q$, and the order initial ideal of the extended right ideal is the
extension of the initial ideal: $\Ann_{K[x,\xi]}\gr Q_K=(\Ann_{k[x,\xi]}\gr Q)\,K[x,\xi]$
(a $k$-basis of $\gr Q$ in each degree stays a $K$-basis; this is the lemma
\texttt{filteredInitialLifting} of the formalization). Consequently
$\Char(Q_K)\subseteq V(P)$ and, since $\Ann\gr Q+(x_1)$ is the unit ideal by
Proposition~\ref{prop:char-avoids}, $\Char(Q_K)\cap\pi^{-1}(H)=\varnothing$.
Alternatively, Sections~\ref{sec:monic}--\ref{sec:char} apply verbatim to $\tilde d\in A_n(K)$.}

Over \AIreplace{$\mathbb C$}{$K$}, we define $W$ as the reduced characteristic variety of $Q$
\[ W=\Char(Q_{\AIreplace{\mathbb C}{K}})_{\mathrm{red}}\]
with $Q_{\AIreplace{\mathbb C}{K}} := Q \otimes_{k_0}\AIreplace{\mathbb C}{K}$.
\AIremove{Under the right-to-left equivalence used above, $(x,\xi)\mapsto(x,-\xi)$
carries $W$ to the characteristic variety of the corresponding left module.} %if gabber statement for right modules i can leave that out
Gabber's involutivity theorem \AIreplace{\mbox{\cite{Gabber1981,Gabber1982}}}{(Theorem~\ref{ithm:gabber}, applied to the right ideal of $A_n(K)$ generated by $\tilde d,x_1^N\tilde d$)} therefore makes
$W$ involutive\AIadd{: its vanishing ideal is $\sqrt{\Ann\gr Q_K}$ by the Nullstellensatz}; see also \cite[Theorem~2.3.1 and Appendix D]{HTT} and
\cite{SinghKumar2014}. \AIreplace{Closedness and fibre-conicality are preserved, as is
containment in $V(P)$, because $P$ is homogeneous of degree $N$ and so
$P(-\xi)=(-1)^NP(\xi)$.}{$W$ is closed, and it is fibre-conical because
$\Ann\gr Q_K$ is homogeneous for the grading by $\xi$-degree ($\gr Q_K$ is a
graded module), so $V(\Ann\gr Q_K)$ is stable under $\xi\mapsto t\xi$.} 
By eq. \eqref{eq:char-contained}\AIadd{ and the base-change remark}, $W \subset V(P)$, by eq.  \eqref{eq:P-e1} $P(1,0,\dots,0) = 1\neq 0$, and by eq. \eqref{eq:char-avoids}, $W\cap \pi^{-1}(H) = \varnothing$. By Theorem \ref{thm:coisotropic-exclusion} $W$ must therefore be empty. 

A variety is empty, iff its reduction is empty, since reducing does not change the underlying point-set, so $\Char(Q_{\AIreplace{\mathbb C}{K}})=\varnothing$.
A nonzero \AIadd{finitely generated module with an exhaustive filtration satisfying $F_{-1}=0$ has nonzero associated graded module, and a nonzero finite module over $K[x,\xi]$ has} \AIremove{coherent filtered module has} nonempty \AIremove{associated-graded} support\AIadd{ (every proper ideal lies in a maximal ideal)}.\AIremove{\texttt{\string\cite\{somewhere\}}}
Consequently $\Char(Q_{\AIreplace{\mathbb C}{K}})=\varnothing$ implies $Q_{\AIreplace{\mathbb C}{K}}=0$. Since field extensions are always faithfully flat \AIremove{\texttt{\string\cite\{somewhere\}}}\AIadd{(a $k$-vector space $V$ with $V\otimes_kK=0$ is zero)}, $\AIreplace{\mathbb C}{K}$ is a faithfully flat module over \AIreplace{$k_0$}{$k$}, and therefore $Q$ vanishes also over \AIreplace{$k_0$}{$k$} \cite[Tag 00HP]{Stacks}
\[Q_{\AIreplace{\mathbb C}{K}} = 0 \Leftrightarrow Q = 0.\]
By eq. \eqref{eq:I-Q}, it follows that 
\[A_n(k_0) = \tilde d A_n(k_0) + x^N \tilde d A_n(k_0)\]
or equivalently, there exist $R,S\in A_n(k_0)$ with
\[1=\tilde dR+x^N\tilde dS.\]
\AIremove{This identity of course extends to the original field $k\supset k_0$.}
Lastly, we can apply the inverse of the automorphism $\Phi_g$, defined in \AIadd{Remark~}\ref{rem:transport}. This results in the final expression
\[ 1=d\bigl(c^{-1}\Phi_g^{-1}(R)\bigr)+F\,d\bigl(c^{-1}\Phi_g^{-1}(S)\bigr)
  \qquad F=\Phi_g^{-1}(\AIreplace{x}{x_1})^N,\]
proving Theorem \ref{thm:main}.
\AIadd{\formalref{\leandecl{Stafford38/FoundationClosure.lean}{37}{Stafford38.FoundationClosure.canonicalSupportVanishingViaGeneralCoisotropic}:
exclusion over $\overline k$ \leandecl{Stafford38/Geometry/GeneralCoisotropicCanonicalAdapter.lean}{107}{Stafford38.Geometry.GeneralCoisotropicCanonicalAdapter.algebraicallyClosedCanonicalSupportVanishing_of_generalCoisotropic},
descent \leandecl{Stafford38/Weyl/FilteredScalarLifting.lean}{453}{Stafford38.Weyl.FilteredScalarLifting.canonicalSupportDescent},
certificate \leandecl{Stafford38/Characteristic/CanonicalCertificate.lean}{30}{Stafford38.CharacteristicCanonicalCertificate.exists_fixedSource_certificate_of_orderCharacteristicSupport_eq_empty}}}
\begin{AIaddition}{Base change of the order initial ideal}
\begin{lemma}\label{lem:initial-base-change}
Let $K/k$ be a field extension, $I=\tilde dA_n(k)+x_1^N\tilde dA_n(k)$ and
$I_K=\tilde dA_n(K)+x_1^N\tilde dA_n(K)$. Then
$\operatorname{in}(I_K)=\operatorname{in}(I)\cdot K[x,\xi]$.
Consequently $\Char(Q_K)$ is the preimage of $\Char(Q)$ under
$\Spec K[x,\xi]\to\Spec k[x,\xi]$. Moreover
$Q=0\iff\operatorname{in}(I)=(1)\iff\operatorname{in}(I_{\Kbar})=(1)$.
\end{lemma}
\begin{proof}
Identify $A_n(K)=A_n(k)\otimes_kK$ through the PBW basis. Then
$F_L(A_n(K))=F_L\otimes_kK$ and $I_K=I\otimes_kK$ as $K$-subspaces. Since
$K$ is flat over $k$,
$(I\otimes_kK)\cap(F_L\otimes_kK)=(I\cap F_L)\otimes_kK$. Hence every
$z\in I_K\cap F_L(A_n(K))$ is a sum $\sum_ic_iw_i$ with $c_i\in K$ and
$w_i\in I\cap F_L$. Its order-$L$ symbol is
$\sum_ic_i\sigma_L(w_i)\in\operatorname{in}(I)K[x,\xi]$. The reverse
inclusion holds because $\sigma_L$ commutes with the scalar extension.

For the last claim: if $\operatorname{in}(I)=(1)$, then $\gr Q=0$. Since
$F_{-1}Q=0$ and the filtration is exhaustive, induction on $m$ gives
$F_mQ=0$ for all $m$, so $Q=0$. The converse is clear. Finally, a proper
ideal of $k[x,\xi]$ has a zero in $\Kbar^{2n}$ (Nullstellensatz), so it
stays proper in $\Kbar[x,\xi]$.
\end{proof}
\end{AIaddition}
\AIadd{\formalref[{stated for the canonical ideal; the argument uses nothing else}]{%
\leandecl{Stafford38/Weyl/FilteredScalarLifting.lean}{384}{Stafford38.Weyl.FilteredScalarLifting.target_orderInitialIdeal_le_map_source_orderInitialIdeal};
\leandecl{Stafford38/Weyl/PresentedScalarExtension.lean}{388}{Stafford38.Weyl.PresentedScalarExtension.presentedWeylScalarExtension_map_orderInitialIdeal_le};
\leandecl{Stafford38/Characteristic/EmptySupportVanishing.lean}{80}{Stafford38.CharacteristicEmptySupportVanishing.orderInitialIdeal_eq_top_iff_rightQuotient_subsingleton};
\leandecl{Stafford38/Characteristic/GeometricSupportDescent.lean}{96}{Stafford38.Characteristic.GeometricSupportDescent.map_scalarPolynomialMap_eq_top_iff}}}
\AIcomment{COMP-15}{\textbf{Information --- proof comparison; no edit requested.} The formal main proof runs the geometry over $\Kbar$ and descends the order initial ideal using Lemma~\ref{lem:initial-base-change}; it does not use the embedding into $\mathbb C$. The rank-zero and scalar cases are handled separately. For the remaining cases, emptiness of characteristic support gives $1\in I$, and the inverse symplectic change returns the fixed-source identity over $k$. \href{https://itpplasma.github.io/stafford38-formal/lean_proof_details.pdf\#nameddest=comparison-COMP-15}{Mathematical proof (COMP-15, p.~9).}}

\AIreplace{Latly}{Lastly}, as Stafford already showed, from $1 = dR+FdS$ follows cyclicity of the torsion right module:
\AIadd{\begin{corollary}\label{cor:cyclic}
Every finitely generated torsion right $A_n(k)$-module is cyclic.
\end{corollary}}
\AIcomment{ALG-06}{\textbf{Proposed change --- author acceptance pending.} State the consequence for \emph{finitely generated}
torsion right modules, as required by Stafford's two-generator theorem.
Write $\phi(r,s)=t_1r+t_2s$. The map $\phi$ is surjective, not an
isomorphism. Correct the final
quotient to $T\cong(A_n\oplus A_n)/M$. The displayed right-module
generator calculation itself is valid.}
\begin{proof}
    Suppose for every $d\neq 0 \in A_n$, there exists $F, R, S \in A_n$, such that $1 = dR + FdS$. Let us denote a torsion module by $T$. \AIreplace{Stafford's theorem (Theorem 3.7 \mbox{\cite[p.~438]{Stafford1978}}) established, that every finitely generated torsion module can be generated by 2 generators,}{Stafford's theorem \cite[Theorem~3.7, p.~438]{Stafford1978}, stated there for left modules and valid for right modules by the left--right symmetry noted on \cite[p.~437]{Stafford1978}, shows that every finitely generated torsion right $A_n$-module is two-generated (the Lean proof below avoids this theorem),} so let us denote the generators of $T$ with $t_1$ and $t_2$. We can define a homomorphism 
    \[\phi:A_n \oplus A_n\rightarrow T \quad \phi(r,s)\mapsto t_1r+t_2s.\]
    This \AIreplace{isomorphism}{map} is surjective, since $t_1, t_2$ generate all of $T$. By the first isomorphism theorem (for $\phi:M\rightarrow N$, $im(\phi) \cong M/ker(\phi)$), it follows
    \[T \cong (A_n\oplus A_n)/M\]
    with $M = ker(\phi)$. \\
    $T$ is torsion and therefore $\forall t\in T: \exists m \in A_n\backslash\{0\}: tm = 0.$ Specifically, let us denote with $d_1$ and $d_2$ the annihilators of $t_1$ and $t_2$, i.e.
    \[t_1d_1 = t_2d_2 = 0.\]
    $A_n$ is a right Ore domain, which means it has the property that for any nonzero $d_1, d_2$ there exists elements $r_1, r_2 \in A_n$ such that 
    \[d_1r_1 = d_2r_2 \neq 0.\]
    Defining therefore $d := d_1r_1 = d_2r_2$ yields us a common annihilator of $t_1$ and $t_2$.\\
    $(d, 0)$ and $(0,d)$ are elements in $M = ker(\phi)$, since $\phi(d, 0) = t_1d = 0$ and $\phi(0, d) = t_2d = 0.$\\
    Because $(0,d), (d, 0) \in M$, also 
    $(d, 0)S + (0, d)R = (dS, dR) = (dS, 1-FdS) =: (a, 1-Fa) \in M$, i.e. $(a, 1-Fa)$ is an element of $M$. \\
    On the other hand, $A_n\oplus A_n$ can be generated by $(a, 1-Fa)$ and $(1, -F)$. We can check this by finding $u,v$ such that for all $(x,y)\in A_n\oplus A_n$ we can write 
    \[(x,y) = (a, 1-Fa)u + (1, -F)v.\]
    Suitable choices are $u = y+Fx$ and $v = x-a(y+Fx)$.\\
    Lastly, since  $A_n\oplus A_n$ can be generated by $(a, 1-Fa)$ and $(1, -F)$ and $(a, 1-Fa)$ is an element of $M$, it follows that $T \cong \AIreplace{/A_n\oplus A_n)/M}{(A_n\oplus A_n)/M}$ can be generated by \AIadd{the image of} $(1, -F)$, i.e. $T$ is cyclic. 
\end{proof}
\AIcomment{COMP-16}{\textbf{Information --- proof comparison; no edit requested.} The formal cyclicity proof uses the identity for each nonzero annihilator and induction on a finite generating set. It does not assume an Ore condition or invoke a separate two-generator theorem. The right-sided placement of the cofactors is essential. \href{https://itpplasma.github.io/stafford38-formal/lean_proof_details.pdf\#nameddest=comparison-COMP-16}{Mathematical proof (COMP-16, p.~10).}}
\AIadd{\formalref{\leandecl{Stafford38/TorsionCyclicity.lean}{161}{Stafford38.TorsionCyclicity.weyl_isCyclic_of_isRightTorsion},
two generators \leandecl{Stafford38/TorsionCyclicity.lean}{27}{Stafford38.TorsionCyclicity.exists_span_singleton_eq_span_pair},
domain \leandecl{Stafford38/Weyl/Domain.lean}{22}{Stafford38.WeylDomain.mul_ne_zero};
Mathlib-only statement \leandecl{CorollaryChallenge.lean}{37}{Stafford38CorollaryChallenge.torsionCyclicStatement}}}

\section{Outlook}
\AIremove{We have proved}
\AIadd{The cofactors $R,S$ are obtained non-constructively: the argument proves
$Q=0$ without producing a certificate. Effective bounds for $\deg_BR$ and
$\deg_BS$ in terms of $\deg_Bd$, and an algorithm in the spirit of the
constructive two-generator methods, remain open. The Global Stafford
extension raises the same questions for rings of differential operators on
smooth affine varieties.}
\AIcomment{EDIT-01}{\textbf{Action required --- unresolved proof or drafting issue.} This sentence and the Outlook section are unfinished.
Please complete the intended outlook or remove this fragment once reviewed.}



\clearpage
\bibliographystyle{amsplain}
\bibliography{references}






























\begin{appendix}


% Evidence: project ledgers, archived outward session messages, model-use
% records, independent reviews, and canonical formal-proof provenance.
\section{AI-assisted research and formalization}\label{app:ai-workflow}

Proof discovery, counterexample searches, review, and formalization was
assisted by agentic AI tools. The work began with
cyclic-extension computations in Macaulay2 and Maple. Attempts to construct
universal cofactors directly produced special cases and reductions but
eventually stalled. The decisive change was to study the canonical quotient
$Q=A_n/(dA_n+x^N dA_n)$: Euler surjectivity and monicity constrained its
characteristic support, while Gabber's involutivity and asymptotic conormal
geometry forced that support to be empty. This turned cofactor selection
into a quotient-vanishing problem. Full Lean formalization was deferred
until the paper argument was complete, although small formal checks
continued throughout. In the final formal development, support avoidance
over the coordinate hyperplane was proved by a localized filtered Koszul
length contradiction.

The workflow separated proposed ideas from verified claims. One controller
maintained the authoritative proof state, while parallel workers explored
distinct mechanisms and returned evidence with an explicit scope.
Independent reviewers checked fixed revisions for quantifiers, sidedness,
circularity, and missing hypotheses. The project protocols were based
on a \texttt{math-frontier} agent skill developed for this purpose: freeze the
claim, identify its first unresolved dependency, and specify how a route
could fail. Conceptual prompts requested a small number of distinct
mechanisms; paper-only tasks sought complete arguments before implementation.
A condensed representative prompt, adapted from the archived protocols, is:
\begin{quote}\small
Work on one precisely stated unresolved claim. Propose at most three
distinct mechanisms, each with an intermediate lemma, a noncircular
justification, and an adversarial test. Preserve the quantifiers and
right-sided cofactors. For a paper task, supply a complete argument; for
formalization, prove the required hypotheses without adding axioms or
using \texttt{sorry}. Return evidence and the first unsupported bridge; do not
promote a conditional result to a theorem.
\end{quote}

Research and review sessions used GPT-5.5, GPT-5.6 Luna and Sol,
GPT-6 Astra, and Claude Fable 5, Haiku 4.5 and Opus 5. The project protocols
assigned routine implementation and replay to Luna and conceptual
consultation to Sol; independent reports also record Opus reviews.
Subsequent manuscript, proof-map, build and release work used Fable 5.1,
Opus 5, Sonnet 5 and GPT-6 Luna. These records establish participation and task roles,
without assigning a particular mathematical idea to one model.

The completed proof uses Lean~4 and Mathlib and pins AlgebraicAnalysis v0.3.0
\cite{Lean4,Mathlib2020,AlgebraicAnalysis030}; the later library archive is v0.3.2~\cite{AlgebraicAnalysis032}.
Formalization distinguished theorems proving required inputs from
conditional theorems merely consuming them. Statement comparison and
independent NanoDa and Lean kernel checks verified the submitted declarations
\cite{Comparator,NanoDa,PalomarDocs}; their proofs use only
\texttt{propext}, \texttt{Classical.choice}, and \texttt{Quot.sound}.
The Stafford38 and Global Stafford headline proofs are registered separately in
Palomar~\cite{Stafford38Palomar2026,GlobalStaffordPalomar2026}; their Lean developments are archived~\cite{Stafford38Zenodo120,GlobalStaffordZenodo106}.

\section{Further work}
\AIcomment{REF-01}{\textbf{Proposed change --- author acceptance pending.} The original appendix's undefined labels
\texttt{prop:symbol} and \texttt{eq:charinP} are visibly replaced below
with the existing order-symbol and characteristic-containment labels.}
\AIreplace{AI also proved some related stuff regarding Right Ore localizations, the corresponding left-hand identity, and apparently the corresponding identity over other rings such as differential rings with laurent series coefficients; I am leaving that out for now.}{The formalization also proves the following consequences, which are not otherwise used here:
the identity in every right Ore localization of $A_n(k)$
\formalref{\leandecl{Stafford38/LocalizationCorollaries.lean}{28}{Stafford38.LocalizationCorollaries.s38_rightOreLocalization}};
the left-handed identity $1=Rd+SdF$ via the formal adjoint
\formalref{\leandecl{Stafford38/LeftHandedCorollary.lean}{20}{Stafford38.LeftHandedCorollary.leftHanded_of_universalStatement}};
the identity for finite-order differential operators on localizations of $k[x_1,\dots,x_n]$, including principal opens, Laurent-polynomial coefficients and the fraction field
\formalref{\leandecl{Stafford38/LocalizedDifferentialCorollaries.lean}{80}{Stafford38.LocalizedDifferentialCorollaries.s38_unconditional_localized_differential}};
and, independently of the geometric argument, an evolutionary version with arbitrary coefficients commuting with one Weyl pair
\formalref{\leandecl{Stafford38/EvolutionaryCorollary.lean}{33}{Stafford38.Evolution.evolutionaryCorollary}}.}
\AIcomment{SCOPE-01}{\textbf{Proposed change --- author acceptance pending.} Review the scope before accepting this paragraph. The coefficients in the stated localization result are Laurent polynomials, not arbitrary Laurent series, and the operators have finite order. Retain further consequences only with their exact hypotheses and proof references; the declaration index is in the technical supplement.}
\subsection{Noncharacteristic H}
An apparently nice \AIremove{(though I do not understand its value)} side theorem that comes from "restriction to the coordinate hyperplane" (after \AIadd{Corollary~}\ref{cor:charP}\AIadd{)}, is that $H = \{x\AIadd{_1}=0\}$ is noncharacteristic\AIadd{; noncharacteristic restriction is what makes the pull-back of $Q$ to $H$ well behaved (Cauchy--Kovalevskaya--Kashiwara)}:



The  nonzero conormal covectors of $H$ are the multiples
$\lambda\,dx$ (or in coordinates $(\lambda, 0, \dots,0)$), with $\lambda\ne0$. Since $P$ is homogeneous and
$P(1,0,\ldots,0)=1$, we have $P(\lambda,0,\ldots,0)=\lambda^N\ne0$, therefore $\lambda\,dx \notin V(P)$. 
Since $\Char(Q)\subseteq V(P)$ by \eqref{eq:char-contained},
$\Char(Q)$ therefore contains no nonzero conormal covector to $H$. More precisely,
the projection
\[
 V(P)\cap T^*\mathbb A^n|_H\longrightarrow T^*H
\]
is finite: on coordinate rings it adjoins $\xi_1$ subject to the monic
equation $P(\xi_1,\ldots,\xi_n)=0$. Its restriction to $\Char(Q)$ is
therefore finite, which is the noncharacteristic condition
\cite[Definition~15.6]{SchnellNotes}.



\begin{definition}[Noncharacteristic, {\cite[Def.~15.6]{SchnellNotes}}]\label{def:noncar}
Let $i\colon H\hookrightarrow X$ be a smooth closed subvariety and $M$ a
coherent $\mathcal D_X$-module. Let
$\varpi\colon T^*X|_H\to T^*H$ be the canonical surjection. Then $H$ is
\emph{noncharacteristic} for $M$ if the restriction
$\varpi|_{\Char(M)\cap T^*X|_H}\colon\Char(M)\cap T^*X|_H\to T^*H$ is a finite
morphism.
\end{definition}
\AIadd{\formalref[{for $X=\mathbb A^n$, $H=\{x_1=0\}$ and a characteristic ideal $J$: the coordinate-ring map $k[T^*H]\to k[x,\xi]/(J+(x_1))$ is finite}]{%
\leandecl{Stafford38/NoncharacteristicHyperplane.lean}{70}{Stafford38.NoncharacteristicHyperplane.IsNoncharacteristic}}}

\begin{proposition}[$H=\{x=0\}$ is noncharacteristic for $Q$]\label{prop:noncar}
With $P$ as in \eqref{eq:P} and $H=\{x_1=0\}$, the projection
$V(P)\cap T^*X|_H\to T^*H$ is finite. Hence $H$ is noncharacteristic for $Q$.
Equivalently, $\Char(Q)\cap T^*_HX$ is contained in the zero section.
\end{proposition}

\begin{proof}
Write $x'=(x_2,\dots,x_n)$, $\xi'=(\xi_2,\dots,\xi_n)$, so that
$T^*X|_H=\Spec k[x',\xi_1,\xi']$ and $\varpi$ forgets $\xi_1$. By
Proposition~\AIremove{\texttt{\string\ref\{prop:symbol\}}}\AIadd{\ref{lem:order-symbol}}, $P$ is monic of degree $N$ in $\xi_1$ with
coefficients in $k[\xi']$; therefore
\[
  k[x',\xi']\longrightarrow k[x',\xi_1,\xi']/(P)
\]
is a finite ring extension, generated by $1,\xi_1,\dots,\xi_1^{\,N-1}$. Thus
$V(P)\cap T^*X|_H\to T^*H$ is finite, and by \AIremove{\texttt{\string\eqref\{eq:charinP\}}}\AIadd{\eqref{eq:char-contained}} so is its
restriction to $\Char(Q)\cap T^*X|_H$. For the last statement, the conormal
space is $T^*_HX=\{x_1=0,\ \xi'=0\}$, on which
$P(\xi_1,0,\dots,0)=\xi_1^{\,N}$; so $V(P)\cap T^*_HX=\{\xi=0\}$ is the zero
section.
\end{proof}
\AIadd{\formalref[{the zero-section statement holds without $x_1=0$: tangential covariables vanishing forces $\xi_1=0$}]{%
\leandecl{Stafford38/NoncharacteristicHypersurface.lean}{129}{Stafford38.NoncharacteristicHypersurface.canonical_principal_hypersurface_finite}, \leandecl{Stafford38/NoncharacteristicHyperplane.lean}{232}{Stafford38.NoncharacteristicHyperplane.canonical_isNoncharacteristic_annihilator},
\leandecl{Stafford38/NoncharacteristicHyperplane.lean}{328}{Stafford38.NoncharacteristicHyperplane.canonicalSupport_conormal_subset_zeroSection}}}
\end{appendix}


\end{document}
